\quellseite{719}

\part{Das Buch Hiob}

\section{Vorrede über das Buch Hiob}

Das Buch Hiob handelt diese Frage: Ob auch den Frommen Unglück von Gott widerfahre? Sie stehet Hiob fest und hält, daß Gott auch die Frommen ohne Ursache allein zu seinem Lob peiniget. Wie Christus Joh. 9, 3. von dem, der blind geboren war, auch zeuget.

Danwider setzen sich seine Freunde und treiben ein groß und lang Geschwätz, wollen Gott Recht erhalten, daß er keinen Frommen strafe; strafe er aber, so müsse derselbe gesündiget haben. Und haben so einen weltlichen und menschlichen Gedanken von Gott und seiner Gerechtigkeit, als wäre er gleich wie Menschen sind, und sein Recht, wie der Welt Recht ist.

Wiewohl auch Hiob, als der in Todes-Nöthen kommt, aus menschlicher Schwachheit zu viel wider Gott redet, und im Leiden sündiget und doch darauf bleibet, er habe solch Leiden nicht verschuldet für andern, wie es denn auch wahr ist. Aber zuletzt urteilet Gott, daß Hiob, indem er wider Gott geredet hat im Leiden, unrecht geredet habe; doch was er wider seine Freunde gehalten hat von seiner Unschuld für dem Leiden, recht geredet habe. Also führet dieses Buches Historia endlich dahin, daß Gott allein gerecht ist, und doch wohl ein Mensch wider den andern gerecht ist auch für Gott.

Es ist aber uns zu Trost geschrieben, daß Gott seine großen Heiligen also lässet straucheln, sonderlich in der Widerwärtigkeit. Denn ehe, daß Hiob in Todes-Angst kömmt, lobet er Gott über dem Raub seiner Güter und Tod seiner Kinder. Aber da ihm der Tod unter Augen gehet, und Gott sich entzeucht, geben seine Worte Anzeige, was für Gedanken ein Mensch habe (er sei wie heilig er wolle) wider Gott, wie ihn dünket, daß Gott nicht Gott, sondern eitler Richter und zorniger Tyrann sei, der mit Gewalt fahre, und frage nach niemands gutem Leben. Dies ist das höchste Stück in diesem Buch. Das verstehen allein die, so auch erfahren und fühlen, was das sei, Gottes Zorn und Urtheil leiden, und seine Gnade verborgen sein.

Sonderlich aber ist hie zu merken, daß der Text klar anzeiget, daß Gott dem Teufel Erlaubniß gebe, auch die Frommen anzufechten, V. 6. Und daß Krieg, Feuer vom Himmel, große Winde und schädliche Wetter vom Teufel aus Verhängniß Gottes verursacht und gemacht werden. Darum soll man wider solchen Feind, der den Menschen zu schaden alle Länder durchsucht, Gott bitten, daß er uns durch seine lieben Engel wider ihn schützen und für dergleichen Jammer behüten wolle. Wie der böse Feind hie selbst zeuget, Gott habe bisher den Hiob verwahret, gleich mit einer Mauer, durch seine Engel. Begehret aber, daß Gott solchen Schutz wegnehmen und ihm erlauben soll. Solches ist nicht also zu verstehen, als habe der Teufel mit Gott geredet, und um Erlaubniß gebeten. Denn von solchem Thun kann man anders nicht reden, denn wie es unter uns Menschen gehet. Darum läßt es die Schrift bei solcher Einfalt bleiben, sonst könnten wir nichts davon verstehen. Daß der Teufel die Welt durchwandert, V. 7., ist wahr, denn er ist an allen Orten, und siehet, wo er könne Schaden thun. Daß er aber für Gott kommt, ist nicht zu verstehen, als käme er in Himmel, sondern auf das einfältigste, daß er ohne Gottes Bewilligung nichts könne ausrichten. Darum heißet Gott fragen oder zu Gott kommen, anders nichts, denn auf Gottes Willen sehen, ob ers verhängen oder zusehen wolle, wie es der Satan im Sinn hat.

\chapter[1]{1}

\section{Vorrede}

Im ersten Kapitel ist eine Historia, wie Hiob mit mancherley Unglück auf einmal umgeben, V. 15., 16., 17., 18., und dennoch Gott darin gelobet, und alles geduldig gelitten habe, V. 27. Uns zum Exempel, daß wir dergleichen, wenn es uns übel gehet, auch thun sollen.

\section{Hiobs Widerwärtigkeit und Geduld}

a Nicht daß er so reich und gewaltig sei gewesen, sondern um seiner Weisheit, Verstand und Gottseligkeit willen ist er herrlicher gehalten, denn andere.
\quellseite{720}
ᵇ Segnen) das ist: Fluchen und lästern.

\section{Gebet}

GOtt, unser treuer lieber Vater, erbarme dich unser, und beschütze uns wider den Teufel durch seine heiligen Engel, daß er nicht auch bei uns Jammer und Noth anrichte, und helfe, daß wir in all unserm Kreuz Ihn auch loben, und das in Geduld tragen und überwinden, Amen.

\chapter[2]{2}

\section{Vorrede}

Geliebte Christen, St. Petrus spricht 1. Epistel am 5., V. 8.: Euer Widersacher, der Teufel, gehet umher wie ein brüllender Löwe, und suchet, welchen er verschlinge. Derwegen wir nicht sicher sein sollen, als wäre der Teufel viele Meilen von uns; sondern sollen wachen und beten. Und St. Paulus 1 Thess. 5, 16. vermahnet: Seid allezeit fröhlich, betet ohne Unterlaß, und seid dankbar in allen Dingen; denn das ist der Wille Gottes in Christo Jesu an euch. Derwegen sollen wir Gott auch unter dem Kreuz loben. Denn was von der Hand Gottes kommt, das ist nicht böse, so wir es mit Glauben empfahen. Denn den Gläubigen muß alles zum Besten dienen, Röm. 8, 28. Ferner, der heilige Prophet Jeremias sagt im 49. Kapitel B. 12.: So spricht der Herr: Siehe, diejenigen, so es nicht verschuldet hatten, den Kelch zu trinken, müssen trinken. Derwegen, ob wir schon unschuldig sein, so sollen wir doch Gott zu ehren, was er nur will, in Geduld leiden, weil der unschuldige Sohn Gottes um unsertwillen den bittern Tod erlitten hat, daß er uns aus der Hölle und vom ewigen Tode erlösete. \unsicher{Erzählter} Sprüche und Lehren erinnert auch das 2. Kapitel.

\section{Hiob am Leibe geplaget, wird betrachtet und besucht}
\quellseite{721}

a Haut für Haut) das ist, für seine Haut läßt er fahren Kinder, Vieh, Gesinde und aller andern Haut.

b (Segne Gott) ja du thust fein, lobest und dienest Gott, und gehest drüber zu Grunde.

\section{Summarium}

Dies ist eine Historia, in welcher wir sehen, wie der Teufel anhält und wo er mit einer Anfechtung nichts kann ausrichten, sich je mehr und mehr mit höhern und schwerern Anfechtungen an dem Menschen versucht, ob er ihn zu Falle, in Ungeduld und Sünde bringen könne. Sie sollten wir von dem frommen Hiob lernen, daß wir, wenn's übel gehet, uns auch trösten, wie er sich tröstet, und des Unfalls nicht beschweren, sondern Gott auch im Leiden loben, der uns unser Lebenlang unzählig viel Gutes, beides an Leib und Seele gethan und aus viel Unglück väterlich geholfen hat.

Und hie zeiget der Text das fürnehmste Stück an, das in diesem Buch gehandelt wird, nämlich, daß Gott sagt, er habe Hiob ohne Ursache also geschlagen, B. 3. Auf daß wir lernen, daß auch die Frommen und Gottesfürchtigen unter das Kreuz müssen, ob sie es gleich nicht wie andere Leute, so ohne Gottesfurcht leben, verdienen.

\section{Votum}

Unser HErr JEsus Christus, der in die Welt kommen ist, daß er des Teufels Werk zerstöre, der wehre des Teufels List und Gewalt, und tröste und erhalte unter allem geistlichen und leiblichen Kreuz alle, die seinen heiligen Namen loben und anrufen, und helfe, daß wir allezeit in Gottesfurcht leben, Amen.

\chapter[3]{3}

\section{Vorrede}

Geliebte Christen! gleich wie der Aramant so hart ist, daß er nicht kann zerschlagen werden: also tragen die Gläubigen das Kreuz nicht allein mit geduldigem, sondern auch mit starkem und unüberwindlichem Herzen und Gemüth; gleichwie aber dennoch der starke Aramant, dem das Eisen und Feuer nichts abbrechen kann, durch frisches und warmes Bocksblut zerbrochen wird: also, wenn die Gläubigen Gottes brennenden Zorn fühlen, und das heiße höllische Feuer ihnen ein wenig ins Herz scheinet und leuchtet, so gerathen sie oft dahin, daß sie wider Gott murren und wollen lieber todt sein, denn leben. Aber also muß unser HErr Gott die großen Heiligen demüthigen, daß sie lernen sich selbst recht erkennen, wie ein elend Gemächt ein Mensch sei, wo er nicht unter dem Schutz und Schirm des Höchsten ist. Es sollen auch die betrübten, angefochtenen Herzen und verzagten Gewissen dies merken, daß sie nicht verzweifeln, wenn sie auch dergleichen Ungeduld und Murren unter dem Kreuz wider Gott fühlen; sondern sollen sich trösten, daß die größten Heiligen sowohl Fleisch und Blut gehabt, und auch, ja so leichtlich gestrauchelt und gefallen sind, als wir; Gott aber habe sich ihrer erbarmet und sie zu Gnaden angenommen; der werde uns auch mit gnädigen Augen ansehen, und uns die Sünde vergeben, um unsers HErrn JEsu Christi willen, und uns erretten und erhalten. Davon lehret auch das 3. Kapitel.
\quellseite{722}

\section{Hiobs Wehklage und Ungeduld}

a Meines Leibes) daraus ich geboren ward, das ist der Mutter Leib.

b Wüste) die mit Bauen umgehen, da zuvor nichts steht.

c Verborgenen) das ist, aus der Erden.

b Bedeckt) was soll der leben, der für Angst nicht weiß, wo aus, wohin, besser todt etc.

\section{Summarium}

Sie gehet an das rechte Leiden Hiobs. Der Teufel durch seine Anfechtung hatte das angerichtet, daß er gesucht hatte, nämlich, daß Hiob nicht mehr spricht, wie zuvor: Der HErr hats gegeben, der HErr hats genommen, der Name des HErrn sei gelobet. Item: Haben wir Gutes empfangen von Gott, warum sollten wir das Böse nicht auch annehmen? Sondern er klaget und schreiet über sein Leiden dermaßen, daß er nicht schweigt und wollte, er wäre nie geboren, B. 2. 3. 4. Also siehet man, wie auch heilige Leute, wenn das Leiden anhält, ungeduldig werden und Gottes Willen nicht können mit Gehorsam ertragen. Denn Fleisch und Blut lässet von seiner Art nicht, auch in den Heiligen. Es murret und klaget, wenn ihm übel ist, wollte des Leidens gern überhoben sein. Solcher Ungehorsam ist eine greuliche verdammliche Sünde. Aber Gott wills gerne vergeben, wenn mans nur für Sünde erkennet und um Gnade und Vergebung bittet.

\section{Votum}

Der barmherzige GOtt vergebe uns unsere Sünde, lindere uns auch die wohlverdienten Strafen, und regiere und stärke uns, daß wir unter dem Kreuz mit Murren und Ungeduld nicht sündigen, sondern allezeit in Lieb und Leid ihn loben und preisen, Amen.

\chapter[4]{4}

\section{Vorrede}

Geliebte Christen! welchen anderer Leute Unglück eine Freude ist, die sind ärger, als die un-
\quellseite{723}
vernünftigen Thiere; denn diese sind traurig, winseln und schreien, wenn sie ihre Gattung oder ihres gleichen in Noth sehen und hören; ja solche sind wahrhaftig Teufelsgenossen, denn der ist ein rechter Schadenfroher und seine Lust und Freude daran hat, wenn es den Menschen übel gehet und sie in Herzensangst, Leibesnoth und Güterverlust oder Ehrenschaden gerathen und stecken. Derwegen soll sich jedermann für solcher Unbarmherzigkeit und teuflischer Bosheit und Sünde fleißig hüten. Denn Gott will auch kurzum nicht, daß wir uns anderer Leute Unfall freuen sollen; sintemal wir alle arme Menschen sind, und kann einem jeden begegnen das, so einem andern widerfahren ist. Dazu sind wir alle Sünder und ist ein jeder werth, daß er stets von Gott heftig gestraft werde; und weil es herzlich und schmerzlich weh thut, wenn eines von denen verhöhnet und geplaget wird, die ihn billig trösten sollen; so soll auch niemand andern ihr Kreuz, Leiden und Elend größer und schwerer machen, oder sie mehr mit Worten und Thaten betrüben, wie dem lieben Hiob seine Freunde gethan haben. Als davon das 4. Kapitel also meldet.

\section{Eliphas fängt an Hiob zu bestrafen}

a Ist das) das ist, da siehet man nun, wie fromm du seist, daß dich Gott so strafet.

b Die Löwen und Löwinnen sind die Reichen und Gewaltigen auf Erden, so die Armen unterdrücken.

c Boten) oder Engeln.

\section{Summarium}

Hie gehet die Disputation an, so die Freunde Hiobs mit ihm in seinem Leiden gehalten und ihm dadurch dasselbige nur größer und schwerer gemacht haben. Des frommen Hiobs Herz stund also: er wollte den Jammer alle gerne leiden, wenn ers verdienet hätte, wie es die Welt mit ihrem gottlosen unbußfertigen Leben verdienet. Denn Gott selbst gibt ihm Zeugniß, im 2. Kapitel B. 3., daß er ihn ohne Ursache habe plagen lassen. Solche Meinung wollten ihm seine Freunde ausreden, und sprechen: Gott sei gerecht, darum thue er niemand unrecht. Weil aber Hiob so heftig gestrafet und mit so viel Jammer überschüttet werde, da müsse folgen, er habe solches wohl verdienet und geschehe ihm nicht unrecht, B. 7. Das heißt denn, das Leiden größer und schwerer machen, wenn das Gewissen dazu kommt, daß es eine verdiente Strafe der Sünden sei, und daß es Gott billig über uns habe kommen lassen.

\section{Votum}

Der heilige Geist regiere uns, daß wir betrübte Leute nicht mehr betrüben, sondern herzlich Mitleiden mit ihnen tragen, sie trösten und ihnen rathen und helfen, auf daß sie, durch seinen Beistand, ihr Kreuz in Geduld tragen und zu Ende bringen, Amen.
\quellseite{724}

\chapter[5]{5}

\section{Vorrede}

Geliebte Christen! es ist wahr, daß Gott der Herr sagt zu seinen Gläubigen Jeremiä im 30. Kapitel V. 11.: Züchtigen will ich dich mit Maße, daß du dich nicht unschuldig haltest. Aber das ist auch wahr, daß der heilige Geist durch Salomon lehret in seinen Sprüchen im 17. Kapitel V. 3., da er spricht: Wie das Feuer Silber und der Ofen Gold, also prüfet der Herr die Herzen. Also ist wahr, daß im Buch der Weisheit im 3. Kapitel V. 5. gesagt wird: Die Gerechten werden ein wenig gestraft, aber viel Gutes wird ihnen widerfahren; denn Gott versucht sie, und findet sie, daß sie sein werth sind. Dagegen wie man das Vieh, welches man bald schlachten will, nicht zu schwerer Arbeit gebraucht, und das nicht abmergelt, daß es die Haut kaum schleppen möge, sondern man mästet es: also läßt Gott die Gottlosen hie eine Zeit machen, was sie nur wollen, und in allem Bösen fortfahren; aber da sie nicht Buße thun, so belohnet ihre Sünde, nach den zeitlichen Strafen, das ewige höllische Feuer. Dieser Lehren soll uns erinnern das 5. Kapitel.

\section{Eliphas fährt in seiner Bestrafung fort}

a Heiligen das ist, zeige mir einen Heiligen, der unschuldig ist geplaget, wie du meinst. Aber die Tollen und Unwitzigen heißet er hie die losen, frechen Leute, die nach Gott nicht fragen; solche verderbet wohl der Zorn und Eifer Gottes.

b Hungrige und Durstige heißt er die Räuber und Tyrannen.

c Nicht gehet das ist, der Mensch verdienet solches mit Sünden, sonst käme es ihm nirgend her.

d Dein Bund das ist, die Steine werden dein Getreide bewahren, weil davon eine Mauer umher gemacht wird.

\section{Summarium}

Das ist gleich die Meinung des vorigen Kapitels; denn Eliphas sagt, die tägliche Erfahrung weise aus, daß die Bösen um ihrer Sünde willen gestraft werden, wiederum, daß Gott den Frommen helfe, B. 2. Darum vermahnet er Hiob, er soll ja nicht gedenken, daß er solchen Jammer nicht habe verdienet, sondern sich frei
\quellseite{725}
für einen Sünder erkennen, so werde Gott wieder mit seinem Trost kommen, und ihm helfen, V. 17. Solches ist der Vernunft nach recht geredet.

Aber Gott führet sein Werk höher, denn es die Vernunft verstehen kann, strafet nicht allein darum, daß man es verdienet habe, sondern daß er prüfen wolle, ob man auch fest am Vertrauen auf seine Hülfe und Gnade halten, und ihm den schuldigen Gehorsam auch im Kreuz und Leiden leisten und seinen Willen geduldiglich tragen wolle.

\section{Votum}

Der treue Gott verleihe, daß wir und alle Kreuzträger unter dem Kreuz seinen Willen erkennen; und helfe, daß wir in der Probe allzeit wohl bestehen, behalten den Glauben und ein gut Gewissen, Amen.

\chapter[6]{6}

\section{Vorrede}

Geliebte Christen! Gleich wie David sein Leid getragen und bitterlich geweinet hat über Saul und Jonathan, da sie durchs Schwert gefallen waren, 2 Sam. 1, 17.; und wie St. Paulus sich des kranken Epaphroditus herzlich annimmt, und seinethalben trauerte, und hoch bekümmert war, Phil. 2, 27.: also sollen wir mit unserm Nächsten in seinem Kreuz, als in Armut, Krankheit und andern Trübsalen, auch herzlich Mitleiden tragen und ihm behülflich und rathsam sein, so fern und viel wir mögen. Denn wer Barmherzigkeit seinem Nächsten weigert, der verlässet des Allmächtigen Furcht, sagt heute Hiob, B. 14. Aber es gehet, wie Sirach sagt und klaget im 37. Kapitel B. 3.: Ach wo kömmt doch das böse Ding her, daß alle Welt so voll Falschheit ist? Wenn es dem Freunde wohlgehet, so freuen sie sich mit ihm; wenn es ihm aber übel gehet, werden sie seine Feinde. Dafür hüte sich jedermann. Darauf höret an das 6. Kapitel.

\section{Hiobs Verantwortung wider Eliphas}

a Das Wild das ist, ihr habt gut trösten, euch mangelt nichts. Man isset nicht ungesalzene, wenn mans besser weiß; aber ich muß wohl selbst dies und das etc.

b Nicht verleugnet das ist, hab ichs doch nicht verdienet, daß ich so geplagt werde, wollte Gott, ich wäre doch todt.

c Verschmachten) das ist, weil meine Freunde jetzt vorüber rauschen, wie ein Wasser, und kennen mich nicht, wird sie auch einmal eine Hitze drücken, so werden sie versiegen und vertrocknen, darum, daß sie mich jetzt verlassen.

d Die Wege Thema) das ist, sie halten mit denen, die mich beraubet haben, wie oben Kapitel 1. steht, gehen denselben recht, und mir unrecht.
\quellseite{726}

\section{Summarium}

Dies ist eine Klagrede des frommen Hiobs, daß er neben allem Unglück, das Gott auf ihn geworfen hatte, auch noch dies leiden und tragen müßte, daß seine Freunde, die ihn billig trösten sollten, mit ihrem Reden nur verzagter und blöder machen, als die neben der Leibesnoth ihm auch das Gewissen beschweren und verzagt machen. Denn wo ein Berg vor zerschlagen und voll Kümmerniß ist, daselbst soll man nicht mehr martern und plagen, sondern trösten und auf Gottes Gnade und Barmherzigkeit weisen. Gleich wie man einen Wagen, der auf die Seiten hängen und fallen will, halten und nicht ganz zu Boden reißen oder stoßen soll. Denn dasselbige geschieht ohne das, wenn er hänget und sinkt, daß er leichtlich gar fällt. Solches thaten die Tröster nicht, die dem Hiob nichts anders konnten fürsprechen, denn daß er ein Sünder wäre, der solche Strafe wohl verdienet hätte, denn Gott thäte niemand unrecht etc. Aber Hiob wehret sich stark und spricht Nein dazu, solches ist nicht wahr; darum wie ihr mich tröstet, werdet ihr in eurer Noth und Jammer wieder getröstet werden.

\section{Votum}

Der barmherzige Gott erleuchte und erweiche unsere Herzen, daß wir auch mit unserm nothleidenden Nächsten Mitleiden haben, und ihm tröstlich und hülflich sein; und steure allen bösen und falschen Herzen, welche die Betrübten mehr betrüben, ängsten, quälen und plagen, Amen.

\chapter[7]{7}

\section{Vorrede}

Geliebte Christen! Es ist wohl wahr, daß viel Elend und Jammer auf Erden ist. Aber es ist alles wie nichts zu rechnen gegen der greulichen Anfechtung, wenn das Herz Gottes feuerbrennenden Zorn fühlet. Denn, wenn dem Menschen sein Herz absagt, so ist ihm immer zu Sinne, Gott habe ihn verlassen und verstoßen, und stehet noch weiß nicht, wie er aus der unaussprechlichen Angst und Noth kommen soll: Wo er sich hinwendet und kehret, stehet er nichts, denn den jämmerlichen Anblick, daß Gott zornig ist, und ihn weder sehen noch hören wolle. Und das ist ein Spiegel, wie es in der Hölle zugehen wird, in der ewigen unaussprechlichen Qual, da das Gewissen in den Gottlosen wird jämmerlich zagen, und der Zorn Gottes über ihnen bleiben; und das wird das rechte Herzeleid über alles Leid sein, und vielmehr brennen, als das höllische Feuer, wenn einem jeglichen sein eigen Herz absagen und ihn überzeugen wird, daß ihm nicht unrecht geschehe. O Herr, treuer Gott, laß uns bedenken den großen Ernst, daß wir klug werden und christlich leben. Das erinnert auch das 7. Kapitel.

\section{Hiobs Klage über der Menschen Elend}

a) Vergeblich, das ist, ich habe Ruhe und der Arbeit ein Ende gesucht, aber das ist umsonst, es bleibet noch immer Unruhe.
\quellseite{727}

\section{Summarium}

Ist eine ernstliche Klag-Rede, daß Gott so viel Jammers und Herzeleid auf die elende menschliche Natur wirft, so sie doch an ihr selbst schwach, gebrechlich, und aller Gefahr und Unglück unterworfen, dazu für dem Tod keinen Augenblick sicher ist, daß wohl Noth wäre, daß Gott schonet und hülfe. Aber Gott (spricht Hiob) legt sich noch selbst wider uns, drängt, drücket und ängstet uns. Darum wäre nichts besser, denn nur die Sünde vergeben und bald gestorben. Solche Worte alle gehen aus der angebornen Ungeduld und Ungehorsam, daß der alte Adam ihm Gottes Willen nicht gefallen lassen, sondern immer andere Mittel und Wege fürschlagen will, denn es Gott ordnet.

\section{Votum}

Gott sei uns allen gnädig, lindere uns die wohlverdienten Strafen, und erhalte uns und alle Kreuzträger in seinem Gehorsam, daß wir nicht aus Ungeduld, mit Worten oder Werken, wider ihn und unsern Nächsten sündigen, Amen.

\chapter[8]{8}

\section{Vorrede}

Geliebte Christen! Von Glück und Unglück urtheilet die Vernunft also: Daß, welchem Menschen es in seinem Leben übel gehet, der müsse einen ungnädigen Gott haben. Und wiederum: Wo alles ist, was das Herz begehret, da wohne Gott. Aber die heil. Schrift zeuget dagegen, daß welchen der Herr lieb hat, daß er denselben züchtige, Ebr. 12, 6. Und so die nur sollten einen gnädigen Gott haben, die in diesem Leben gute geruhige Tage haben, so würde daraus folgen, daß die heiligen Patriarchen, Propheten, Apostel und Märtyrer alle von Gott verstoßen wären. Daher stellet auch der Herr Christus das Urtheil von dem Blindgebornen, Joh. 9, 3., und sagt: Es hat weder dieser gesündiget, noch seine Eltern; sondern, daß die Werke Gottes offenbar würden an ihm. Dieß soll wohl betrachtet werden, daß niemand Gott in sein heimlich Gericht falle und seinen Nächsten richte und verdamme; wie mancher sich seines Nächsten Unglückes freuet, und weiß nicht, daß das seine so nahe für der Thür ist, und unterdeß grünet und blühet. Darauf höret auch davon an das 8. Kapitel.

\section{Hiob wird von Bildad der Heuchelei beschuldiget}
\quellseite{728}

\section{Summarium}

Hie ist die andere Disputation, welche Bildad mit Hiob hat; der stehet auch auf dem, Gott richte nicht unrecht; derhalben was Hiob leide, das leide er seiner Sünde halben, V. 3. Er soll aber frömmer werden, denn er bisher gewesen sei, so werde ihm Gott gnädig sein und aus dieser seiner Noth helfen, V. 5. Solche Meinung beweiset er mit der Erfahrung, daß Gott je und je die Frommen erhalten, und die Bösen verworfen und gestraft habe. Also muß die Vernunft schließen, welche nicht mehr von Gottes Willen und Werken weiß, denn im Gesetze offenbar ist; aber der arme Hiob bedarf eines andern Trostes und einer höhern Weisheit in seiner Noth und Anfechtung.

\section{Votum}

Gott, der HErr, vergebe es allen, die betrübte Kreuzträger mehr ängsten; und erleuchte sie, daß sie Buße thun und von solcher Sünde ablassen, und tröste alle Christen in ihrer Noth, und helfe uns auch gnädiglich im Leben und Sterben, Amen.

\chapter[9]{9}

\section{Vorrede}

Geliebte Christen! Wenn ein Kreuz hin ist, so denkt der Mensch, er sei nun gar hindurch. Aber in einem Nu findet sich ein anderes, daß Gott wieder vorüber gehet, ehe wirs gedenken, und verbirget sein Angesicht, läßt uns sinken, jedoch nicht gar ertrinken, und stellet sich, als kenne er uns nicht mehr. Aber das ist gewiß, wenn uns dünket zur Zeit der Noth, Gott sei am weitesten, so ist er uns am nächsten. Diese seine wunderliche Gegenwärtigkeit hat er in der Sonne uns vorgebildet. Wenn uns dünket mitten im Winter, sie sei am weitesten von uns, so ist sie etliche tausend Meilen Weges in ihrem Zirkel am Himmel niedriger, und dem Erdboden näher, denn mitten im Sommer, wenn um Johannis-Tag die größte Hitze ist. Gleich so ist der allmächtige Gott uns viel näher zur Zeit der Noth, denn wenn es gehet, wie wir selbst wollen. Ob aber wohl Fleisch und Blut solches nicht begreift, so findet sich es doch also endlich, wenn wir hindurch sind, und überwunden haben. Darauf höret nun an das 9. Kapitel.

\section{Hiobs Vertheidigung wider Bildad}

a Orion ist das helle Gestirn gegen Mittag, das die Bauern den Jacobstab heißen. Die Glucke oder die Henne sind die sieben kleine Gestirne.

b Die stolzen Junker, die sich auf ihre Macht verlassen und jedermann helfen können.
\quellseite{729}

c Scilicet, ne redeant. D. i. daß sie nicht wiederkehren.

d Kleider) das ist, meine Tugend.

\section{Summarium}

Ist eine Antwort auf Bildads Rede, in welcher Hiob noch immer darauf beharret, er sei nicht gottlos gewesen, habe auch solchen Jammer nicht wohl verdient. Daß aber Gott solches über ihn verhänget, das thue er als ein Gott, der Herr über alles ist, und mag mit jedermann machen, wie er will, dem niemand nichts einreden noch abgewinnen könne, V. 4. Solche Antwort ist etwas hart, aber der Schmerz und die große Noth und Anfechtung treiben Hiob, der von Herzen betrübt war, daß er also redete; sintemal er bekennen mußte, daß andere Leute so gottlose und böse gewesen, dennoch von Gott so hart nicht sind gestraft worden, als er, der doch gottesfürchtig und fromm gewesen war.

\section{Votum}

Gott, der heilige Geist, regiere und stärke uns alle, daß wir auch in der Frömmigkeit beständig bleiben, und was uns darüber widerfähret, daß wir es in christlicher Geduld tragen, und Gott allezeit loben und preisen, bis zu unserm seligen Ende, Amen.

\chapter[10]{10}

\section{Vorrede}

Geliebte Christen! In den Gläubigen, ob sie schon Gottes heilige und auserwählte Kinder heißen, streitet auch allwegen hienieden auf Erden in ihren Herzen der Geist wider das Fleisch, und der Glaube wider den Unglauben. Denn jetzt haben sie gewiß und glauben stark, Gott sei ihnen gnädig und barmherzig, und lache sie freundlich an; bald gerathen sie wieder in schwere Gedanken: Gott zürne mit ihnen, sei ihnen spinnefeind, habe ihnen den Rücken zugekehrt, seine Barmherzigkeit ihnen entzogen, und wolle ihr ganz und gar keine Gnade haben. Dieser Wechsel in den Christgläubigen höret nicht auf bis in die Grube. Derhalben soll niemand gedenken, daß die Christen allzeit stark im Glauben, brünstig in der Liebe und bissig im Gebet sein können und bleiben. Und fromme Christen sollen sich desto fleißiger nach dem Ziel strecken und Gott bitten, daß er ihnen Victoria und Sieg durch den Glauben verleihen und erhalten wolle, Amen. Darauf höret auch an das 10. Kapitel.

\section{Hiob klaget über seinen Jammer}
\quellseite{730}

a Um und um) nichts ist an mir, das du nicht gemacht hast, oder nicht dein sei. Noch verwirfst du mich, als hätte mich ein anderer gemacht, der dein Kind wäre, sogar nimmst du dich deines Eigenen nicht an.

b Odem) das ist, mein Leben, das der Odem anzeigt.

\section{Summarium}

Ist eine Klagrede, in welcher Hiob des fürnehmlich sich beklagt, er müsse leiden und wisse doch nicht warum. Denn er sei nicht gottlos gewesen, und wisse, daß Gott ein Wohlgefallen an den Frommen habe.

Darum klaget er, daß Gott, von welchem wir das Leben und alles haben, sich mit solcher Macht wider einen armen Menschen legt, V. 14.; und wünschet abermals, daß er nicht geboren wäre, V. 18. Also siehet man, wie Natur und Vernunft aus großer Ungeduld wider Gott murret, und des Trostes vergißet, welchen Gott allwege hinter das Leiden stellet und verberget, ob er gleich die Menschen läßt plagen, daß er doch wiederum wolle helfen, wo man sich Hülfe zu ihm versiehet, und bei ihm suchet.

\section{Votum}

Der Gott alles Trostes und der Geduld behüte uns und alle christlichen Kreuzträger für Ungeduld und allen andern Sünden; und erhalte sie und uns in wahrem Trost und allen christlichen Tugenden, bis an unser Ende, Amen.

\chapter[11]{11}

\section{Vorrede}

Geliebte Christen! Den betrübten Kreuzträgern soll man freundlich, tröstlich und herzlich zusprechen, und sie vermahnen, daß sie das liebe Kreuz, so ihnen Gott der Herr auferlegt, geduldig tragen und sich unter die gewaltige Hand Gottes, damit er unterdrückt, aber auch wieder aufhilft, demütigen sollen. Man soll sie trösten, daß solch ihr Kreuz und Leiden zu Gottes Ehre und ihrer Seelen Seligkeit durch Christum gereichen werde, und daß ihnen Gott zu seiner Zeit werde abhelfen. Und wenn sie ohne alles Kreuz wären, so möchten sie zu Weltkindern und mit der Welt verdammt werden. Auch werde ihnen solch Kreuz dazu dienen, daß sie denn ihren treuen Gott, wenn er ihnen wiederum aus allen ihren Anfechtungen herrlich helfen wird, werden von Grund ihres Herzens danken, loben und preisen, daß er sie so väterlich gezüchtiget und durch das liebe Kreuz für vielen Sünden behütet hat; und werden mit allen
\quellseite{731}
lieben Heiligen das Freuden-Liedlein aus dem 118. Psalm V. 18. singen: Der Herr züchtiget mich wohl, aber er gibt mich dem Tode nicht. Also soll man mit betrübten Christen handeln, und die nicht mehr betrüben und anschnauzen, wie Zophar dem leidenden Hiob thut, in dem 11. Kapitel.

\section{Hiob wird von Zophar der Heuchelei beschuldiget}

a Wib) das ist, frei und seines Willens.

\section{Summarium}

Ist die dritte Disputation, Zophars von Naema; der stehet auch darauf, wäre Hiob fromm gewesen, als er sich macht und hören läßt, so hätte ihn Gott nicht so gestraft. Weil ihn aber Gott also straft, so sei es ein gewisses Zeugniß, daß ers verdienet habe, und sei gottlos gewesen, V. 13. Denn Vernunft kann höher nicht kommen, denn daß sie von Gott urtheilet, wie ein Mensch von dem andern pflegt zu urtheilen. Da gehets also zu, wer unrecht thut, der muß entgelten. Wer fromm ist, der genießet. Aber Gottes Urtheil und Gericht ist viel höher. Der läßt den frommen Hiob greulicher strafen, denn alle Gottlosen. Denn er gedenket erstlich, er wolle ihm aus solchem Jammer wieder helfen, und ihn desselbigen reichlich ergetzen. Zum andern will er sehen, ob Hiob auch gerecht und schlecht bleiben, in Leiden nicht ungeduldig werden und von Gott nicht abweichen wolle. Solcher Versuchung bedarf Hiob wohl sammt allen Heiligen.

\section{Gebet}

Der allmächtige, gerechte, gütige GOtt schaffe seinen Willen auch an uns, und helfe, daß wir den erkennen, und alles, was er uns zu leiden zuschickt, mit Geduld annehmen und tragen, und im Glauben auch allezeit bis zum seligen Ende verharren, Amen.

\chapter[12]{12}

\section{Vorrede}

Geliebte Christen! Gott, der HErr, hält lieb und werth die Frommen und Gottesfürchtigen. Daher gebeut er ernstlich in seinem Wort, daß ihm dieselben niemand soll betrüben, beschädigen noch beleidigen. Denn also sagt er im 105. Psalm V. 15.: Tastet meine Gesalbten nicht an, und thut meinen Propheten kein Leid. Also spricht er auch Zach. 2, 8.: Wer euch antastet, der tastet meinen Augapfel an. Derwegen so nehme ihn nur keiner in Sinn, daß es ihm für wohl und ungestraft werde hinaus gehen, der sich an den Christen und Auserwählten Gottes vergreift, und ihnen entweder mit Worten oder Werken zu nahe und wider ist. Denn die bedrängten Christen sagen und klagen Gott ihre Noth; welch ihr Gebet und Seufzen ein unleid-
\quellseite{732}
lich Geschrei ist für den Ohren Gottes, daß ein Mensch lieber die ganze Welt zum Feinde haben sollte, denn ein einziges Christen-Gebet, das mit Ernst wider ihn wird gesprochen. Denn es sind eitel Donnerschläge, was sie in solcher Noth reden, die gewißlich ihre Feinde treffen. Dies soll auch erinnern das 12. Kapitel.

\section[Hiob fängt an sich wider Zophar und die]{Hiob fängt an sich wider Zophar und die andern Freunde zu verantworten}

a Id est, glimmend Docht.

b Das ist: Ihr saget: Weisheit sei bei den Großvätern; ich sage aber, sie sei bei GOtt, welcher gleich aller Könige, Priester, Richter, Gewalt, Kunst, Heiligkeit zu nichts machet.

\section{Summarium}

Ist eine Antwort auf Zophars Rede, in welcher Hiob sagt, er wisse es sehr wohl, daß GOtt allmächtig sei und mit uns thun könne, was er will. Denn solches sehe man an allen Creaturen, V. 6. Aber nichts desto weniger (spricht er) sehe und erfahre ichs, daß auch den Frommen übel und den Bösen wohl geht. Auf die Frage (spricht er) sollt ihr mir eine richtige tröstliche Antwort geben, so wollte ich euch für weise Leute halten, daß andere weiß ich fürwahr. In Summa, so Hiob auf andere Leute stehet, ist er frömmer, denn sie, und dünket ihn zuviel, daß es ihm sollte übler gehen, denn ihnen. Und fehlet ihm durchaus an dem, daß er nicht von Herzen beten noch sprechen kann: Dein Wille geschehe auf Erden, wie im Himmel.

\section{Votum}

Der treue, gütige GOtt helfe und verleihe, daß wir uns seinen väterlichen Willen allezeit wohlgefallen lassen, auf daß wir in christlicher Geduld alles, was er uns auflegt, willig tragen, dankbarlich rühmen und ritterlich überwinden, Amen.

\chapter[13]{13}

\section{Vorrede}

Geliebte Christen! Es ist ja ein unbeständig Ding um das menschliche Leben. Denn gleich wie der Wind im Herbst die kahlen verdorreten Blätter dahin führet; also ists auch bald um einen Menschen geschehen. Desgleichen, wenn der Böttcher ein Faß pichtet, so gehet ein Rauch auf, als wollte er die Sonne verdunkeln; aber gar bald, wenn ein kleines Windlein drein gehet, so wird er vertrieben: also auch wir, wenn uns der kalte Wind des Todes anhaucht, so ist es

44
\quellseite{733}
Alles aus. Und an dem Fluß Hypane findet man Würmlein, welche, wenn die Sonne aufgehet, lebendig werden, so bald sie aber untergehet, so sterben sie: also kurz ist auch unser Leben für Gott, wie eine Nachtwache, Ps. 90, 4., daß wir auch nicht wohl eines Tages gewiß versichert sein. Derwegen soll niemand in Tag hinein sicher leben, und meinen, er habe mit dem Tode einen Verstand gemacht, Jes. 28, 15. Sondern ein jedes soll in Gottesfurcht sich üben und sein Leben desto besser anwenden, weil es so kurz und flüchtig ist, und soll sich befleißigen, daß er einen guten Namen hinter ihm lasse, daß er Gott und den Menschen nützlich in seinem Leben allhie gewesen sei. Darauf höret ferner an das 13. Kapitel.

\section{Hiob fährt in seiner Beantwortung fort}

\section{Summarium}

Dies ist eine Widerlegung deß, daß Zophar dem Hiob hatte fürgeworfen, und ist dieß ungefähr die Meinung: Was ihr mir saget, spricht er, weiß ich fürwahr, darum will ichs nicht annehmen, daß ihr von Gott saget: Er strafe allein die, so es mit ihren Sünden verdient haben. Ihr seid unnütze, verdrießliche Aerzte, V. 4. Man soll die Leute nicht also trösten, so zuvor Gottes Strafe auf ihnen haben. Denn also macht man das Böse nur ärger und des Jammers mehr. Ja, spricht er weiter, ihr selbst, so ihr nichts anders von Gott wisset, denn das allein, daß er ein gerechter Gott sei, werdet ihr solchen Gott nicht können ertragen. Denn seine Gerechtigkeit wird euch zu schwer sein. Ich bin ein armer geplagter Mann, aber die Hoffnung habe ich, Gott werde mein Heil sein und mich nicht halten für unrecht, wie ihr mich beschuldiget.

Auf diese Rede fähret Hiob an zu beten, daß Gott mit seiner Strafe ablassen und nicht so hart gegen ihm handeln wolle, V. 20. Denn er sei ein dürres Blatt; habe er in seiner Jugend gesündiget, soll ers ihm vergeben und gnädig sein, also siehet man, daß in Anfechtung das Gebet

* Beißen) das ist, was soll ich mich viel fasten und mit wehe thun, so ich doch sterben muß und halt mich nicht. Item, meine Seele in die Hände legen, das ist, viel wagen und in Gefahr gehen.
\quellseite{734}
der größte Trost ist, an welches alle Heiligen sich gehalten haben. Und alsdann sind sie auf dem rechten Wege, da sie nicht können unrecht thun, wenn sie sich schuldig geben und Gnade begehren. Merke aber, wie auch der fromme Hiob seiner Jugend halben sich schuldig gibt, V. 26., wie David auch. Denn die Jugend ist zu schwach und das Fleisch zu stark, und kann derselben sich selbst nicht regieren. Darum wer durch solchen Strudel will, muß nicht sicher sein und fleißig beten.

\section{Votum}

O Gott regiere und stärke uns, daß wir alle die Sicherheit ja fliehen fleißig, in rechtem Glauben beten und in wahrem Troste alle Widerwärtigkeit ertragen, bis zum seligen Ende, Amen.

\chapter[14]{14}

\section{Vorrede}

Geliebte Christen! So bald wir mit Weinen, Schmerzen und Traurigkeit auf dieser Welt geboren werden, gehet eines jeglichen Menschen Angst und Noth bald an, und, je länger er lebet, je mehr Trübsal und Elend ihm unter Augen gehet; und wächst das Unglück und Kreuz mit ihm auf täglich, daß wohl recht gesagt wird: Was ist der Mensch jetzt nach dem Fall? Nichts denn Unglück hie und überall, ist doch nichts mehr das elende Leben, als daß wir stets im Tode kleben, so lang es währet. Wohl dem geschieht, der sanft im Frieden hindurch bricht. Dieses alles sollen die Christen in Geduld tragen und in Gottes Wort und bei Gott, durch Gebet, Hülfe und Trost suchen und bedenken, wenn es hoch kommt, daß gleich der Leib darüber zu Trümmern gehet und zu Asche und Pulver wird, so soll er doch herfür am jüngsten Tage und ähnlich werden dem herrlichen verklärten Leibe unsers Herrn Jesu Christ. Dieser Trost macht, daß alles Trauern verschwindet und lindert alle Schmerzen des Leibes, daß es die Christen ertragen können und überwinden. Dieses erinnert auch das 14. Kapitel.

\section{Von des menschlichen Lebens Nichtigkeit}

*Hoffnung) das ist, für dem Tode hat er keine Hoffnung in diesem Leben.

\section{Summarium}

Dies ist eine Klage, in welcher Hiob klaget über den Jammer, welchen alle Menschen auf

44*
\quellseite{735}
Erden haben, daß ihr Leben kurz und voll Unruhe ist, V. 1., und der Mensch nichts gewisser für ihm hat, denn das Sterben, und hernach im Tode liegen ohne alle Hoffnung, daß er wieder möge zu diesem Leben kommen, weil Himmel und Erde stehet.

Darnach bittet Hiob, daß Gott an jenem Tage sein in Gnaden gedenken und alle Sünden vergeben wolle, V. 13. Denn mit diesem Leben hie auf Erden sei es nichts, denn Mühe und Arbeit. Des Menschen Hoffnung, spricht er, ist verloren, das ist, weil wir hie auf Erden sind, können wir nicht hoffen, daß es mit uns besser werden soll. Es muß allerlei erlitten und erstritten und endlich gestorben sein. Gott sei uns gnädig an jenem Tage, wenn Himmel und Erde vergangen sind, Amen.

Dies Zeugniß von der Auferstehung der Todten und einem andern Leben nach diesem Leben, soll man hie fleißig merken. Denn in der ganzen Schrift des Alten Testaments kaum ein klarers ist, und ist das allererste, daß er des ewigen Lebens vermagenden gedenkt, daß er Gottes Gnade begehrt und anruft, und nicht, wie zuvor, auf seine eigene Frömmigkeit stehet, noch derselben sich tröstet. Denn er findet, daß er zu gering zu solcher Wohlthat ist, und dennoch stehet das ganze Papstthum darauf, wir können es mit eigenen Werken und Verdienst erlangen. Aber es scheinet, daß sie in der Schule nicht sind gewesen, da Hiob in ist.

\section{Votum}

Helfe Gott, der allein helfen kann, daß wir den jämmerlichen Zustand unsers Lebens erkennen, den täglich fleißig betrachten, zu täglicher wahrer Buße gebrauchen, in wahrem Glauben endlich einschlafen und am jüngsten Tage zum ewigen Leben mit Freuden auferstehen, Amen.

\chapter[15]{15}

\section{Vorrede}

Geliebte Christen! Ach schrecklich ist es, daß der Mensch ein Greuel und schnöde ist. Denn der Mensch heißt billig ein Greuel darum, daß er wider alle Gebote Gottes sündiget; schnöde und untüchtig aber darum, daß er das Gute nicht thut, so ihm Gott geboten hat. Schrecklich ist es auch, daß der Mensch wegen solcher Sünden ein Kind des Zorns, ein Gefangener des Teufels, ein Gefäß der Unehre und unter dem Fluch und Gottes Zorn ist. Und es sei einer so weise, verständig, hoch und ansehnlich, als er immer wolle, so sind alle seine herrlichen Tugenden, wie ein güldener Harnisch über einem aussätzigen Körper. Sollten doch alle Haare auf unserm Haupte und alle Blutstropfen im Leibe sich darüber entsetzen, daß die Sünde so ein greulicher Greuel ist, und so grausamen Schaden den armen Menschen zugefüget hat.

Und daraus sollen wir lernen Gott fürchten, die Sünde meiden und eine herzliche Abscheu, Furcht, Entsetzung und Erschreckung dafür haben unser Lebenlang. Ach das helfe uns allen der treue Gott, hochgelobt in Ewigkeit. Ferner höret an das 15. Kapitel.

\section{Eliphas will aus Hiob einen Heuchler machen}

a Geringe), das ist, meinst du, daß Gott die Sünder tröste und seinen Trost so geringe hinwerfe, du mußt zuvor fromm werden.
\quellseite{736}

b Brüstet sich scilicet contra Deum, sicut bos pingvis, et crassus, pugnat. D. i. nämlich gegen Gott kämpft er, wie ein fetter und wilder Stier.

\section{Summarium}

Sie fängt Eliphas zum andern mal an, wider den armen Hiob zu disputiren, und spricht: Bist du ohne Sünde, so bist du mehr und höher, denn nie kein Mensch gewesen ist. Denn keiner unter allen Heiligen Gottes ist ohne Tadel, B. 15. Du aber, spricht er, bist ein Heuchler, läßt dich wohl dünken, du seiest fromm, aber die Strafe weiset anders aus. Denn wie Gott allen Gottlosen dräuet, eben also gehests auch dir.

Also siehet man für und für in diesem Buch, daß die Vernunft von Gott nicht anders kann denken noch urtheilen, denn nach dem Gesetze, welches sagt: Gott wolle die Bösen strafen und den Frommen gnädig sein. Wer aber für Gott böse oder fromm sei, und warum Gott auch die Frommen unter das Kreuz werfe, da weiß die Vernunft nichts davon. Darum will Hiob solche Meinung ihm nicht gefallen lassen, wie folget.

\section{Votum}

GOtt der HErr, als ein lieber Vater, stärke und tröste alle betrübte Kreuzträger, und helfe, daß sie der Menschen Spott und des Teufels Anfechtung allezeit mit Geduld überwinden, und regiere uns, daß wir jedermann tröstlich und förderlich sein und bleiben, Amen.

\chapter[16]{16}

\section{Vorrede}

Geliebte Christen! Der liebe Hiob ist unter seinem schweren langwierigen Kreuz mit eitel Todesgedanken umgegangen, und hat nicht hoffen können, daß er daraus werde errettet werden können, dieweil er weder Maß noch Ziel, weder Mittel noch Wege hat mögen ersehen, wie und wodurch ihm geholfen möchte werden; sondern hat nicht anders gewußt, er müßte schon sterben. Also ist allen betrübten und trostlosen Herzen zu Sinn, daß sie gedenken: Nun ists aus. Und wenn sie keine menschliche Hülfe ersehen können, bestimmen sie ihnen selbst die Zeit, da sie nun gewiß sterben und von der Welt hinweg müssen. Dazu kommen denn auch die Gedanken, die der Satan eingibt, dagegen ihnen das Herz im Leibe erbebet und alle ihre Gebeine erzittern. In solcher Angst und Anfechtung sollen ihnen fromme Christen einbilden die Sprüche heiliger Schrift, die uns überzeugen, daß Gott, der HErr, viel reichlicher und überschwänglicher thun kann über alles, das wir bitten und verstehen, nach der Kraft, die da in uns wirket, als St. Paulus lehret Ephes. 3, 20.; und das David im 4. Psalm B. 4. sagt: Erkennet doch, daß der HErr seine Heiligen, das ist, Gläubigen, wunderlich führet. Dazu höret an das 16. Kapitel.

\section{Hiob bezeuget seine Unschuld}
\quellseite{737}

a Zeugen) das ist ihr Befehl wider mich.

b Horn) das ist, meine Gewalt, Macht und Herrschaft, und worauf ich mich verließ.

\section{Summarium}

Ist eine Antwort auf Eliphas Rede, in welcher Hiob nicht sagt, daß er allein gerecht und ohne Sünde sei. Denn seine Sünde hat er mehr denn einmal bekennet und Gott um Vergebung gebeten. Aber das sagt er, sie seien verdrießliche leidige Tröster. Denn mit solchem Trösten machen sie ihm nur des Jammers mehr, V. 1. Fähet derhalben an und machet eine neue Klage von seinem großen Leiden und Schmerzen, darinnen er ist, daß Gott ihn mit so mächtiger Hand zu Boden geschlagen habe, V. 2., und die bösen Leute ihm solches Unglück mit ihrem leidigen Disputiren noch begehren schwerer zu machen; Gott im Himmel weiß meine Unschuld. Aber es will mich gegen diese Heuchler nichts helfen, darum begehre ich zu sterben.

\section{Votum}

GOtt im Himmel, der gerecht ist, und Barmherzigkeit erfordert, steure allen, welche die Betrübten mehr ängsten und betrüben; und verleihe allen Kreuzträgern Trost und Stärke, daß sie ihr Kreuz in Geduld zu Ende bringen; und regiere uns, daß wir in Heiligkeit und Gerechtigkeit leben, Amen.

\chapter[17]{17}

\section{Vorrede}

Bei diesem Kapitel sollen wir merken das Stück, nämlich, wie sich die Frommen sollen, so in langwieriger Krankheit liegen und immer nach dem Tode seufzen und wünschen. Erstlich sollen sie ja eigentlich und gewiß das fassen, daß sie Gott in ihrer großen Angst und Nöthen durch sein Wort und heiligen Geist also herrlich stärken, tröstigen, gründen und erhalten wolle, und also mit der Anfechtung schicken, daß sie es ertragen und ausstehen werden, wie 1 Cor. 10, 13. geschrieben stehet: Gott aber ist getreu, der euch nicht läßt versuchen über euer Vermögen, sondern machet, daß die Versuchung so ein Ende gewinne, daß ihrs könnet ertragen. Und in Klagliedern Jeremiä am 3. Capitel V. 31.: Der HErr verstößet nicht ewiglich, sondern er betrübet wohl und erbarmet sich wieder, nach seiner großen Güte. Denn er nicht von Herzen die Menschen plaget und betrübet. Auch sollen sie sich trösten, wenn dermaleinst das selige Stündlein kommen wird, daß sie aus der Welt werden scheiden sollen, daß sie denn sanft werden dahin fahren und süß im HErrn Christo einschlafen, und daß das letzte Wort unsers HErrn Jesu Christi, so er mit seinem himmlischen Vater am Kreuz geredet hat, auch ihr
\quellseite{738}
letztes Wort sein wird, wo nicht auf ihrer Zunge, doch in ihrem Herzen, nämlich: Vater, in deine Hände befehle ich meinen Geist, Luc. 23, 46. Zum andern sollen sie sich herzlich freuen auf das ewige Leben und auf die himmlische Freude, so kein Auge gesehen, kein Ohr gehört und in keines Menschen Herz kommen ist, Esaia 64, 4. 1 Cor. 2, 9. Und sollen wissen, daß ihre Seele ins Paradies wird aufgenommen werden und bei dem Herrn Christo sein, und daß sie sanft und lieblich ruhen werden in ihren Gräbern, bis die letzte Posaune erschalle: Stehet auf ihr Todten; und daß alsdenn sie dünken wird, als hätten sie kaum eine Viertel-Stunde geschlafen. Zum dritten sollen sie ihnen auch gegen das gegenwärtige Elend und Trübsal dieser Welt die Freude des ewigen Lebens wohl einbilden, daß allda ein Augenblick wird besser sein, denn die hundert tausend Jahre, daß Gott alsdenn die Thränen wird abwischen von allen Angesichtern, Off. Joh. 21, 4. Daß sie die Herrlichkeit des Herrn Christi sehen werden, die er gehabt, ehe der Welt Grund gelegt ward, Joh. 17, 24. Und daß vor ihm Freude die Fülle ist und lieblich Wesen zu seiner Rechten ewiglich, Ps. 16, 11. Und daß Gott der Herr alsdenn seine Gläubige sättigen wird mit langem Leben und ihnen weisen sein Heil, Ps. 91, 16. Das ist, seinen einigen lieben Sohn, in seiner Glorie, in seiner Majestät und Herrlichkeit.

\section{Hiob erzählt sein Elend und Unschuld noch weiter}

\section{Summarium}

Ist eine Klagrede des armen Hiobs, daß er elend, mühselig und von Herzen betrübt sei, und dennoch ein gut Gewissen habe seines Lebens halben, daß er nicht gottlos sei gewesen. Solches, spricht er, können diese meine Tröster nicht verstehen, B. 10. Es thut aber den Frommen sehr wehe, wo es ihnen also begegnet, wie mir geschund. Aber nichts desto weniger werden sie ihnen solches lassen ein Exempel sein, und darum von der Gerechtigkeit nicht abweichen, ob sie gleich auch etwas leiden müssen.

Daß er aber klaget, er müsse in die Hölle fahren, B. 13. 16., hat nicht die Meinung, wie wir im Deutschen verstehen, als müsse er verdammt sein; sondern in die Hölle fahren, heißt sterben, und im Tode oder unter der Erde liegen, da wir alle hingehören. Aber doch durch Christum und Gottes Barmherzigkeit die Hoffnung haben, daß wir auch mit dem Leib wieder auferstehen und ewig leben werden.

\section{Votum}

Der große starke Gott verleihe großen Trost allen, die groß Kreuz tragen, und erhalte sie in wahrem Glauben mit Leib und Seele zur großen Freude des ewigen Lebens; und aber behüte er nach seinem gnädigen Willen für Krankheit und anderer großen Noth und Elend, Amen.

\chapter[18]{18}

\section{Vorrede}

Dies Kapitel lehret, wie Gott die Gottlosen mit allerlei Plage zu strafen pflege. Derwegen, auf daß wir uns Gottes Zorn und Strafe mit unsern Sünden nicht zuziehen, so wollen wir betrachten, wie wir uns für dem listigen Anlauf des Satans hüten sollen, auf daß er uns nicht in Sünde falle und verführe. Dazu aber gehöret endlich, daß man Gottes Wort fleißig höre, lese, betrachte und also täglich mit umgehe. Zum
\quellseite{739}
andern sollen wir anhalten mit dem Gebet, ernstlich und emsig, und sollen mit Mund und Herzen bitten: Ach Herr, führe uns nicht in Versuchung, wehre du dem leidigen Teufel, erhalte du uns, daß er uns nicht in Sünde und Schande leite und stürze. Zum dritten soll ein jedes wohl bedenken den ernsten Zorn Gottes, und wie schrecklich er den Sündern mit allerlei Plage und Strafe an Leib und Seele dräuet, 3 Mos. 26, 14., 5 Mos. 28, 15. Zum vierten sollen wir auch oft an die Exempel gedenken derer, die von Gott greulich sind gestraft worden, und jämmerlich dahin gegangen, von wegen ihrer großen und schweren Sünden. Zum fünften, wer der Sünde entgehen will, der muß auch die Oerter, Personen und Gelegenheit meiden, dadurch ihm der böse Geist nachstellet. Das sind die Mittel, dadurch uns Gott der Herr wider die List und Gewalt des schädlichen Bölands, des Teufels, bewahret. Es muß aber auch gekämpfet sein wider die Sünden, und muß ihren bösen Reizen widerstanden sein.

\section{Bildad mischet Hiob unter die Gottlosen}

a Versetzt) das ist, Gott wird mit dir nicht anders machen, denn mit allen andern, und seine Weise nicht lassen um deinetwillen.

b Fürst) das ist, die Macht und Gewalt des Todes. Also auch König des Schreckens ist die Gewalt des Schreckens, daß er muß unten liegen und nicht entrinnen kann.

c Wurzeln) Wurzel heißt er alles, was in der Erde gepflanzet ist. Ernte alles, was oben auswächset, es sei Korn, Oel, Wein etc.

\section{Summarium}

Ist eine Disputation des Bildads wider Hiob, in welcher er erzählet, wie Gott die Gottlosen pflege zu strafen mit allerlei Plage. Will also anzeigen, weil eben solche Plagen dem Hiob auf dem Nacken liegen, so sei er auch gottlos und leide solches alles um seiner Sünden willen.

\section{Votum}

Wir sagen Gott Lob und Dank, der uns diese Mittel in seinem Wort offenbaret hat; der verleihe, daß wir sie auch allezeit gebrauchen, in seiner Furcht leben, das Gute thun und das Böse lassen, und zeitlichen und ewigen Strafen entfliehen, Amen.

\chapter[19]{19}

\section{Vorrede}

Geliebte Christen! Dieweil es einem wohl gehet, so hat er viel, die sich seiner annehmen; so bald sich aber das Glück wendet, so fliehen sie davon, wie die Schwalben und etliche andere Vögel Winterszeit. Wenig sind, die in der Noth und Anfechtung bestehen. Unserm Herrn Christo selbst ist das begegnet. Denn da es an die rechte
\quellseite{740}
Noch ließen ihn seine Jünger allein in der Feinde Gewalt. St. Paulus beklagt sich auch, da er für Nero gestellet war, daß ihn seine Freunde, auch Demas, verlassen haben, 2 Tim. 4, 10. Das thut trefflich weh, sonderlich wenn wir uns unserer Freunde viel getröstet und viel auf sie gebauet haben. Aber es läßt uns Gottsolches auch zum Guten geschehen. Denn wir geben etwa den Menschen zu viel zu, so müssen wir erfahren mit unserm Schaden, wie wankelmüthig sie sind, und lernen, daß wir uns Gottes allein trösten sollen, und ihn für unsern Felsen, Schloß, Schild, Schutz und Schirm halten. Wahr ist’s, man muß mit einander haushalten und einander trauen, wie man sagt: Wenn alle Tage auch ein Mord geschähe; aber trauen soll man einem Menschen als einem Menschen. Dies soll auch erinnern das 19. Kapitel.

\section[Hiob redet von seinem Elend, und seinem Trost]{Hiob redet von seinem Elend, und seinem Trost von der Auferstehung}

a Ehre, Krone, Hoffnung ist alles geredet vom zeitlichen Leben in guter Ruhe.

b Satt werden) das ist, könnt nicht aufhören mich zu beißen und zu fressen.

c Erlöser, Retter) Vindex, Quia Christus vindicatos, contra Homicidam nostrum Diabolum. D. i. Rächer, weil Christus uns rächet an unserm Mörder, dem Teufel.

\section{Summarium}

In die Antwort auf Bildads Rede, in welcher Hiob wiederum klaget, daß seine Freunde mit solchen Reden ihm das Leiden nur schwerer machen. Gott (spricht er V. 6.) brauchet solche Macht wider mich, das muß ich leiden und dazu still schweigen. Denn was kann ich wider Gott
\quellseite{741}
ausrichten? Hebt also wieder an und klaget seine Noth. Am Ende aber tröstet er sich über die Maßen sein: wenn schon hie auf Erden muß gelitten und gestorben sein, so weiß ich (spricht er B. 25.) doch, daß mein Erlöser lebt, der wird mich an jenem Tage aus der Erde und dem Tod, zu einem andern bessern und ewigen Leben auferwecken. Auf solchen Trost wollten ihn seine Freunde nicht weisen, darum tröstet er sich selbst also. Uns zum Exempel, daß wirs ihm auch nachthun sollen. Eben auf solche Weise vermahnet Paulus Timotheum, 2 Tim. 2, 8.: Halt im Gedächtnis Jesum Christum, der auferstanden ist von den Todten, aus dem Samen David, nach meinem Evangelio, über welchem ich mich leide bis an die Bande. Denn weil es doch beschlossen ist, daß die Gottseligen dem Exempel Christi, des Sohnes Gottes, nach, hie auf Erden ihr Kreuz tragen und sich leiden müssen, ist dies der einige Trost, daß wir mit Christo wieder vom Tod auferstehen und ewig leben sollen.

\section[Notum]{\unsicher{Notum}}

Lob, Ehr, Preis und Dank sei GOtt für diesen Trost; der verleihe, daß er uns und allen unter dem Kreuz ein beständig Labsal und rechte Erquickung sei; und helfe, daß wir christlich leben und selig sterben, und mit Freuden am jüngsten Tage auferstehen, Amen.

\chapter[20]{20}

\section{Vorrede}

Geliebte Christen! Die Strafe stimmet gemeiniglich mit der Sünde überein, wie das Buch der Weisheit im 11. Kapitel V. 17. sagt: Womit jemand sündigt, damit wird er auch geplagt. Also, Antiochus Epiphanes hat viele fromme Herzen betrübet und beleidiget, darum mußte er endlich selbst für Trauern sterben, 1 Macc. 6, 13. Also, Attila, der Hunnen König, hatte viel Bluts vergossen, darum fraß ihn das strenge Gericht Gottes, daß er in seinem eigenen Blute erstickte. Also auch, es trägt sich oft zu im Ehestande, daß mancher ein böses Weib hat, manche wiederum einen bösen unartigen Mann kriegt, daß ist beiden eine Strafe der Sünden, die sie zuvor, ehe sie in Ehestand getreten sind, begangen haben. Eben also gehts den Geizigen, wenn sie den Armen ihr Gut auffressen und ihren Schweiß und Blut aussaugen, so beraubet sie Gott auch wiederum des Ihren, Jer. 30, 16. Derwegen hüte sich ein jeder für Sünden, daß ihm auf sein Sündendienst nicht Sündenlohn folge. Dies lehret auch das 20. Kapitel.

\section{Zophars Rede von der Gottlosen unbeständigem Glück}

a Wenn er Bosheit ansähet, hat er Wollust und Ruhe; aber es wird nicht währen, wird bald bitter schmecken.

b Saugen) das ist, er wird tödtliches Herzleid und Jammer leiden, und alles Guten beraubet werden.
\quellseite{742}

c Ausgeblasen) das ist, ein Feuer von Gott angezündet, nicht durch Menschen ausgeblasen.

\section{Summarium}

Ist eine Disputation des Zophars wider Hiob. Dem Gottlosen (spricht er) gehet es wohl, aber nicht lange, B. 5. Ursache, wenn sie niemand mit dem Herzen helfen und alles mit ihrem Gebet zu sich reißen, kanns Gott in die Länge nicht leiden, sondern macht sie wieder zu Bettlern, B. 19. Also, spricht er, gebets jetzt dir auch; wärst du fromm gewesen, es sollte dir wohl länger sein wohlgegangen; weil du aber gottlos bist, bleibet das Sprichwort wahr: Des Gottlosen Ruhm stehet nicht lang, und die Freude des Heuchlers währet ein Augenblick. Solche Einrede nimmt Zophar aus der Erfahrung. Denn gewiß ist es, daß Gott die Geizigen wie einen vollen Schwamm wiederum ausdrücken und sie leeren wolle, wie das Sprichwort heißt: Uebel gewonnen, übel zerronnen. Aber die Welt glaubets nicht, bis sie es erfähret.

\section{Votum}

GOtt verleihe, daß alle Geizige und andere Sünder sich bekehren; steure allen, die wider ihn und die Liebe ihres Nächsten handeln, und helfe uns, daß wir der Gerechtigkeit nachjagen und andern gerne und treulich dienen, Amen.

\chapter[21]{21}

\section{Vorrede}

Geliebte Christen! In diesem Leben gehet es einem frommen Christen so bald übel, als einem gottlosen und bösen Menschen. Derwegen ist es ein großer epicurischer Irrthum, daß man schließen will: Siehe, der lebet, wie er auf Rosen ging, derwegen muß er gewiß einen gnädigen Gott haben. Denn, wenn Gott das Böse allwege alhie zeitlich strafe, so bedürfte er des jüngsten Gerichts nicht. Dieweil es aber die Gottlosen oft alhie auf Erden am besten haben, und die Gläubigen dagegen viel Böses leiden müssen; so können vernünftige Leute daraus abnehmen und verstehen, daß ein letztes Gericht und nach diesem Leben eine Belohnung und Strafe sein müsse. Dazu grünen eine Zeitlang wohl die Gottes-Verächter in diesem Leben und stehen in voller Blüthe, aber endlich gehen sie oft auch allhier unversehens zu Grund und Boden. Die Christen aber, wie elend und voller Jammers sie sind, so haben sie doch vor sich die herrliche Erlösung und das ewige Leben. Das soll die Gottlosen schrecken, daß sie Buße thun, und das soll die Gläubigen trösten, daß sie ihr Kreuz in Geduld tragen. Davon lehret auch Hiob im 21. Capitel.

\section[Hiob widerleget die Rede Zophars von dem Glück]{Hiob widerleget die Rede Zophars von dem Glück und der Strafe der Gottlosen}

a Augenblick) das ist, sie leben bis an den Tod wohl und da ist um einen bösen Augenblick mit ihnen zu thun, so sind sie hindurch. Ich aber muß so lange Zeit Schrecken und Unglück leiden.
\quellseite{743}

b Vergelten) das ist, wer kanns urtheilen, was ihm zu vergelten sei, ohne Gott allein?

c Haufen) das ist, es ist ihm ein Grab bereit, unter andern Gräbern.

d Der Schlamm des Bachs) das ist die Wollust im Fleisch und guten Leben.

\section{Summarium}

Ist eine Antwort auf Zophars Urtheil, daß es unrecht sei. Ursache, die Gottlosen haben hie auf Erden in diesem zeitlichen Leben einen Vortheil, daß es ihnen lange wohl gehet, V. 7. Die rechte Strafe aber, die sie verdienen (spricht Hiob), ist nicht das zeitliche Unglück, sondern der Tod und die ewige Verdammniß, V. 17. Darum darf man nach dem, wie es ihnen hie auf Erden gehet (es sei böse oder gut), nicht urtheilen, wie ihr thut, wenn mirs übel gehet, so sprecht ihr, ich sei gottlos gewesen. Ich aber sehe wohl, daß das Widerspiel wahr ist, nämlich, daß die Gottlosen nicht so hart gestraft werden, als ich. Aber ihre verdiente Strafe gehöret in ein anderes Leben, davon ihr nichts wisset. Dies ist wiederum ein helles Zeugniß vom ewigen Leben und Auferstehung der Todten.

\section[Notum]{\unsicher{Notum}}

Der treue GOtt helfe, daß wir von allen Dingen nach seinem heiligen Wort urtheilen lernen! Erhalte alle Betrübte unter ihrem Kreuz zur ewigen Freude, und bestätige uns in der Frömmigkeit, zu seines heiligen Namens Lob, Ehre und Preis, Amen.

\chapter[22]{22}

\section{Vorrede}

Geliebte Christen! Wenn wir uns wegen anderer Leute Unglück freuen, aus Neid, Nachgier und anderer fleischlicher Anfechtung, so ist es unrecht und verboten. So es aber geschieht um Gottes und des Nächsten Willen, so ist es recht und zugelassen. Denn viel Heilige Gottes, wie die Schrift bezeuget, sich nicht allein gefreuet haben, wenn es den Feinden Gottes übel gegangen, sondern sie haben auch gebeten, daß er sie strafen oder gar ausrotten wolle, 2 Mos. 15, 4. 1 Sam. 25, 39. Spr. Sal. 11, 5. Das dienet aber also zur Ehre Gottes, so jedermann siehet, daß Gott zu seinen Dingen schauet, daß er seine Gewalt noch nicht von Händen gegeben, daß er seine Feinde wohl dämmen und seine Gerechtigkeit erzeigen könne. Es dienet dem Nächsten und den Einfältigen zu gut, die sich an der Glückseligkeit des Gottlosen heftig ärgern, ungeschickte Reden vermögen fahren lassen und ihrem bösen Exempel ansehen zu folgen, oder sonst ansehen zu schlottern, daß sie sich eines andern und bessern bedenken. Es dienet auch den bösen Leuten, die gestraft werden, zu gut, daß sie nur menschliche Vernunft haben, daß sie aufhören zu sündigen, da sie nicht mehr und härter sollen gestraft werden. Dies erinnert auch das 22. Kapitel.
\quellseite{744}

\section{Eliphas ziehet Hiobs Frömmigkeit abermal in Zweifel}

a Finsterniß heißt Trübsal und Unglück, wiederum Licht heißt Glück und Heil.

\section{Summarium}

Ist die dritte Disputation, in welcher Eliphas eben die vorige Meinung wiederholet, daß Gott den Frommen gnädig sein, sie behüten und ihnen helfen wolle, wiederum die Gottlosen strafen. Weil nun die Strafe auf dem Hiob öffentlich liegt, so soll ers als verdienten Lohn annehmen, V. 6., seine Sünde bekennen und sich bessern, so werde ihm geholfen werden.

Auf diese Weise muß die Vernunft von Gott urtheilen und reden, anders kann sie nicht, und ist zwar nicht ungerecht geredet; doch so fern, daß man beides bekennet, daß wir nicht recht können fromm sein, und daß Gott wolle barmherzig sein und Sünde vergeben, wenn mans begehrt und hoffet, durch Christum.

\section{Votum}

Unser treuer Heiland JEsus Christus, Gottes und Mariä Sohn, helfe, daß wir unserer Gebrechlichkeit uns täglich erinnern, in wahrem Glauben uns seiner wider unsere Sünde trösten, daß wir uns auch ferner für allen Sünden hüten, und durch ihn endlich selig werden, Amen.

\chapter[23]{23}

\section{Vorrede}

Geliebte Christen! Daß man unter dem Kreuz geduldig sei, wider Gott nicht murre und seiner gnädigen väterlichen und der nöthigen Hülfe
\quellseite{745}
erwarte, so soll ein jedes zweierlei betrachten, glauben und hoffen: Erstlich soll ein jedes herzlich bekennen, daß es sein Kreuz mit seinen Sünden wohl verdienet habe, wie Judith lehret Kapitel 8, V. 22. und saget: Wir sollen bekennen, daß es eine Strafe ist von Gott, viel geringer, denn unsere Sünden sind. Ja wir sollen bekennen, daß wir auch die größten zeitlichen, und darnach auch die ewige Strafe verdienet hätten, wenn Gott mit uns ins Gericht gehen wollte, und mit uns handeln nach unserm Verdienst. Wie wir denn solches alles bekennen müssen, wir haben auch so heilig gelebt, als wir immer vermochten. Fürs andere soll ein jedes sich herzlich trösten, daß Gott wieder helfen wolle und uns endlich von allem Uebel erlösen, wie er selber zusagt Jes. 54, 7.: Ich habe dich einen kleinen Augenblick verlassen, aber mit großer Barmherzigkeit will ich dich sammeln. Ich habe mein Angesicht im Augenblick des Zorns ein wenig von dir verborgen, aber mit ewiger Gnade will ich mich dein erbarmen. Dieses hat auch Gott der heilige Geist den königlichen Propheten David gelehret, derwegen er also im 30. Psalm V. 6. sagt: Des Herrn Zorn währet ein Augenblick, und er hat Lust zum Leben, den Abend lang wäret das Weinen, aber des Morgens die Freude. Dies soll uns erinnern das 23. Kapitel.

\section[Hiob beruft sich um seines guten Gewissens willen]{Hiob beruft sich um seines guten Gewissens willen auf Gottes Richterstuhl}

a Einig also Gal. 3, 20. Gott ist einig. Des Einigen aber ist kein Mittler.

\section{Summarium}

Hie ist eine Antwort auf Eliphas Rede, in welcher Hiob auf diesen zweien Puncten beharret. Erstlich, er sei fromm und habe solche Strafe nicht verdienet, V. 6. 11. Zum andern, daß er aber solche Strafe tragen müßte, da könnte er nicht vor sein, Gott sei allmächtig, der mache es mit den Menschen, wie er wolle.

Solche Gedanken zeiget nicht allein Ungeduld an, sondern ist auch Gott in dem zunahend, als handelte Gott mehr nach Macht, denn nach Recht und Billigkeit.

Darum fehlet es dem frommen Hiob in seiner Anfechtung an diesen zweien Stücken. Das erste, daß er die sündhafte Natur nicht eigentlich erkennt. Denn obgleich andere Leute dem äußerlichen Leben nach ärger sind, denn Hiob, so ist doch Hiob dem Urtheil Gottes damit nicht entbrochen. Denn das sündhafte Fleisch und Blut ist eben sowohl in ihm, als in andern Leuten. Das andere, daß er sich zu Gott nicht versiehet, daß er werde wieder helfen und gnädig sein.

\section{Gebet}

Der treue GOtt erleuchte unsern Verstand, und erweiche unsere Herzen, daß wir unter dem Kreuz unsere Sünden erkennen, uns demüthigen und seiner Hülfe uns trösten und freuen, bis zur seligen Erlösung; auch gläubig und fromm sein und bleiben allzeit, Amen.

\chapter[24]{24}

\section{Vorrede}

Geliebte Christen! Dieweil Gott unser trot mit der ewigen wohlverdienten Strafe, um seines lieben Sohns, unsers Herrn Jesu Christi willen, verschonen will, sollen wir uns nicht weigern, allhier ein wenig Kreuz zu leiden, sondern sollen mit dem Propheten Micha sagen, Mich. 7, 9.: Ich will des Herrn Zorn tragen, denn ich habe wider ihn gesündiget; und sollen gedenken, daß es besser sei, hie leiden, denn in jenem Leben, wie St. Augustinus betet: O Herr baue hie, schlage hie und schone dort. Wie auch Nicephorus vom Kaiser Mauritius schreibt, daß
\quellseite{746}
er solches von Gott gebeten habe, und da ihm vor seinen Augen der Tyrann Phocas sein Gemahl und Kinder umbringen ließ, und er nun auch herhalten sollte, hat er gesprochen: Herr Gott, du bist gerecht, und dein Gericht ist auch recht, Ps. 119, 137. Und derwegen sagt St. Ambrosius: Das ist eine herrliche That, daß du in Schmach Gottes Urtheil lobest, in Krankheiten seinen Willen preisest, und in Armuth seine Gerechtigkeit rühmest. Ferner höret an das 24. Kapitel.

\section[Von Gottes verborgener, doch gerechter Regierung, der Frommen]{Von Gottes verborgener, doch gerechter Regierung, der Frommen Elend und der Gottlosen Glück betreffend}

a) Die Zeiten) weil Gott die Bösen so lässet machen, wie sie wollen, so scheinet es, als wüste er nichts darum. Weil ihr denn saget, er strafe die Bösen und nicht die Frommen, so müsset ihr zugeben, daß er nicht wisse und die ihn kennen, auch nicht wissen, zu welcher Zeit er strafen werde, wie ihr euch rühmet zu wissen.

b) Das Wild) die freien, frechen Leute und Tyrannen.

c) Bauet) das ist, die der Hurerey nachgehen, bringen ihr Gut um und lassen ungebauet.

d) Auf ihr Thun) das ist, daß sie nicht ein Aufruhr wider ihn machen, dämpfet er sie immerdar, und muß also Sicherheit mit List suchen, aber es währet nicht.

\section{Summarium}

Dies ist eine Entschuldigung des armen Hiobs, in welcher er anzeigt, er sei kein Tyrann, B. 2., kein Mörder, B. 14., kein Ehebrecher, noch dergleichen Sünder gewesen, B. 15. Welche, wenn es ihnen gleich übel abgeht, haben sie es wohl verdienet, B. 20. Beruhet also immer darauf, er habe solche Strafe nicht verdienet. Das ist nun wahr, wenn man auf der Menschen Urtheil sehen will. Aber Gottes Urtheil ist höher, da sind wir allezeit bußfällig, ob uns gleich die Menschen nicht können beschuldigen.
\quellseite{747}

\section{Votum}

Der barmherzige GOtt strafe uns nicht, wie wirs verdienet haben, und helfe, daß wir das Kreuz in Geduld tragen, seine Werke allzeit rühmen und preisen, und uns, wie Hiob gethan, für allerlei Sünden fleißig hüten, Amen.

\chapter[25]{25}

\section{Vorrede}

Geliebte Christen! Ob sich schon die Christen hüten vor äußerlicher Sünde, Schande und Laster: dennoch, damit sie nicht möchten gedenken, sie wären ganz und gar ohne Sünden, und also von sich selbst gerecht, so schicket ihnen Gott, der HErr, Kreuz und Leiden zu, welches Bestrafungen oder Belohnungen der Sünden sein, Röm. 6, 23. Damit sie ihre Erbschuld und Erbsünde erkennen, darin sie empfangen und geboren und von des wegen auch von Natur Kinder des Zorns Gottes sind. Wie davon Gott, der HErr, sagt Jerem. am 30. Kapitel V. 11.: Züchtigen will ich dich mit Maße, daß du dich nicht unschuldig hältst. Derwegen sollen wir abermal unter dem Joch des Kreuzes, und unter der Last der Anfechtungen unsere Seele mit Geduld fassen, und sollen Gott recht geben, wenn er uns züchtiget und mit dem Kreuze hart drücket, und mit den Anfechtungen peinlich ängstet, und sollen mit dem lieben David bekennen, und sprechen Ps. 51, 6.: An dir, HErr, allein hab ich gesündiget, und übel vor dir gethan, auf daß du Recht behaltest in deinen Worten, und rein bleibest, wenn du gerichtet wirst. Dieses sollen uns erinnern das 25. und 26. Kapitel.

\section{Vor GOtt sind alle Menschen Sünder}

a Wer soll dir thun? Gott ist allmächtig und kann wohl steuern den Großen, wenn du nur fromm wärest. Und du meinest, er wisse es nicht, wie du jetzt neulich gesagt hast.

\section{Summarium}

Ist eben einerlei Meinung auf die Einrede, die oben Bildad that, Gott ist allein gerecht und rein, für ihm kann sich niemand seiner Reinigkeit noch Gerechtigkeit rühmen. Nein, spricht Hiob, ich bin gerecht, denn ich habe nicht so gottlos gelebet, wie andere. Daß ich aber also leiden muß, geschieht darum, daß Gott thut, was er will, das muß ich geschehen lassen, leiden, und dazu still schweigen. In Summa, es fehlet dem frommen Hiob in seiner Anfechtung an dem, daß er sich ansiehet, wie er gegen andere Leute sei, und vergißet daneben, wie er gegen Gott sei, daß er solchen Gehorsam nie habe geleistet, wie ihn Gott gefordert. Soll derhalben sich schuldig geben und nicht sagen: Gott thut mit mir, was er will. Sondern er thut mir, wie ich verdienet habe, und soll mich noch ewig verdammen, wenn er meinem Verdienst nach und nicht nach seiner Gnade mit mir wollte umgehen. Aber die arme schwache Natur kann in der Anfechtung solches nicht bedenken. Bedenket sie es aber recht, so muß sie verzweifeln, wo Gott nicht hilft.

\section{Votum}

GOtt, der heilige Geist, der in alle Wahrheit leitet, lehre und regiere uns, daß wir unser Kreuz recht ansehen und verstehen, und christlich tragen, und behüte uns und alle betrübte Kreuzträger für Verzweiflung und allen andern Sünden, wider GOtt und wider die Menschen, Amen.

\chapter[26]{26}

\section{Hiob preiset GOttes Majestät viel herrlicher, als Bildad}

a Riesen die großen Wallfische, welche bedeuten, die großen Tyrannen auf Erden.
\quellseite{748}

\chapter[27]{27}

\section{Vorrede}

Geliebte Christen! Der gerechte und starke Gott dräuet Malachiä am 3. Kapitel V. 5.: Ich will zu euch kommen, und euch strafen, und will ein schneller Zeuge sein wider die, so mich nicht fürchten. Und St. Paulus vermahnet Röm. 11, 20. und sagt: Sei nicht stolz, sondern fürchte dich. Hat Gott der natürlichen Zweige nicht verschonet, daß er vielleicht dein auch nicht verschone. Durch diese Sprüche werden wir vermahnet, daß wir uns für Sünden hüten sollen, auf daß wir mit unsern Sünden Gott nicht verursachen, daß er uns strafen müsse, und daß mit Gottes Strafe unser eigen Gewissen begangener Sünde wegen uns auch nicht nage und plage. Denn wie ein gut unverletzt Gewissen ist der edelste beste Schatz und Kleinod auf Erden und ein stetiges Paradies und Wohlleben: also ist wiederum keine greulichere, größere, schmerzlichere und ängstlichere Plage, als ein böses Gewissen, ja es ist die Hölle selber. Dieses soll uns erinnern das 27. Kapitel.

\section[Hiob hält, seine Unschuld zu retten, den Frommen]{Hiob hält, seine Unschuld zu retten, den Frommen und Heuchler gegen einander}

a Heuchler heißen in diesem Buch allenthalben einen falschen Menschen, wie sie alle sind, für Gott ohne Glauben.

b Weinen) sie werden froh werden, daß er todt ist.

\section{Summarium}

Nachdem die Freunde Hiob ausgeredet und ihn immer beschuldiget haben, er habe solche Strafe verdienet, darum soll er sich nicht für fromm halten, sondern vielmehr seine Sünde bekennen und Gott um Gnade bitten. Darauf will Hiob wider sie schließen und sagt: Sie thun ihm unrecht; er habe solche Strafe nicht verdienet, wie es die Gottlosen verdienen, welche Gott mit mancherlei Plagen pflegt zu strafen,

45
\quellseite{749}
N. 12. Weil nun solche Strafen, wie Hiob nach der Länge erzählet, über die Gottlosen sollen kommen, so soll hie billig jedermann sich für dem gottlosen Leben hüten und fromm sein. Denn obgleich das Kreuz und Unglück auch über fromme Leute kömmt, ists doch träglicher, denn ein böses Gewissen samt der Strafe haben.

\section{Votum}

Regiere und behüte uns der treue GOtt, daß wir nicht muthwillig, wissentlich und vorsätzlich sündigen, und da er uns nach seiner Gerechtigkeit begangener Sünden wegen strafte, so verleihe er uns gnädiglich Trost und Freude, unter aller Strafe bis zum seligen Ende alles Uebels, Amen.

\chapter[28]{28}

\section{Vorrede}

Geliebte Christen! Daß Gottes Gnade, Treue und Liebe gegen uns groß sei, kann man daran auch erkennen, daß er uns seine Weisheit durch sein Wort und Geist offenbaret und den theuren Schatz mittheilet. Gedenket aber, wenn ein Fürst einem seiner Unterthanen seine Rathschläge offenbarte, er würde es traun groß und hoch achten; also wenn ein König oder Herzog einen Ring von seinem Finger zöge, darin ein Türkis, Sapphir oder Smaragd versetzt wäre und den einem unter uns gebe, wie herrlich und hoch würden wir den halten? Nun stellet uns Gott in seinem Wort seinen Willen und die rechte Weisheit für, gegen welcher alles Gold, Silber, Edelgestein auf Erden nichts ist. Wer danket aber Gott für solchen großen Schatz? Wer trachtet darnach so fleißig, als wir sollten? Diese große Undankbarkeit läßt uns erkennen und läßt uns allen Fleiß anwenden, daß wir wahre Weisheit erlangen, auch mehren und behalten. Welches geschehen wird, so wir ernstlich Gottes Wort, dadurch sie gepflanzet wird, lernen und betrachten, und da wir Gott, der sie allein gibt, darum bitten werden. Denn also sagt der heilige Apostel Jacobus im 1. Kapitel seiner Epistel, V. 5.: Welchem Weisheit mangelt, der bitte sie von Gott, der da gibt einfältiglich jedermann, und rückets niemand auf, so wird sie ihm gegeben werden. Darauf höret auch an das 28. Kapitel.

\section{Lob der Weisheit an GOtt und Menschen}

a Finstern) das ist, man gräbet zuletzt so tief, daß man findet, das verborgen liegt in Finsterniß der Erde.

b Stolzen Kinder) das sind junge Löwen.
\quellseite{750}

\section{Summarium}

Ist eine feine Lehre, daß nichts so tief unter der Erde und im Meer verborgen sei, man könne ihm nachgraben und es finden, aber die rechte Weisheit könne man nirgend finden, denn allein im Wort GOttes. Dasselbige lehret, wie die rechte Weisheit nichts anders sei, denn GOtt fürchten und das Böse meiden. Solche Weisheit (will nun Hiob zu seinen Freunden sagen) habt und wisset ihr nicht; ich aber weiß und habe sie, darum dürft und könnt ihr mich nicht lehren.

\section[Notum]{\unsicher{Notum}}

GOtt, der du bist der Brunnquell aller Weisheit, gebe uns allen auch durch sein Wort und Geist die rechte Weisheit, und helfe, daß wir hier Sünder uns unser Lebenlang hüten, und in seiner Furcht gehorsamlich allzeit leben, bis zur seligen Heimfahrt aus diesem elenden Leben, Amen.

\chapter[29]{29}

\section{Vorrede}

Geliebte Christen! Der Mann Gottes, Moses, vermahnet alle Welt im 13. Kapitel B. 4. seines 5. Buchs und sagt: Ihr sollt dem HErrn, eurem GOtt, folgen, und ihn fürchten, und seine Gebote halten, und seiner Stimme gehorchen, und ihm dienen, und ihm anhangen. Dazu setzet dies St. Paulus 2 Tim. 3, 12.: Die gottseelig leben wollen in Christo Jesu, müssen Verfolgung leiden. Denn der HErr Christus lehret selbst also Matth. 10, 12.: Wer nicht sein Kreuz auf sich nimmt und folget mir nach, der ist mein nicht werth. Da hören wir, daß wir alle gegen GOtt und gegen der Welt christlich und ehrbar leben sollen, und da das liebe Kreuz uns dennoch darauf folget (wie denn alle Christträger oder gläubige Christen Kreuzträger sein müssen), daß wirs in Geduld auf uns nehmen und tragen sollen; in Betrachtung auch dieses, daß der Sohn GOttes, unser HErr Jesus Christus, den größten vollkommenen Gehorsam seinem himmlischen Vater in Thun und Leiden, nach seinem heiligen höchsten Willen geleistet, und hat dennoch das allergrößte, schwerste und längste Kreuz an Leib und Seele tragen müssen, und hat solches gethan in tieffer Demuth und unaussprechlicher Geduld. Dies soll auch erinnern das 29. Kapitel.

\section{Hiob erzählt seine vorige Glückseligkeit}

a In Butter) da ich alles übrig genug hatte, alles fett und vollauf.

b Bogen) das ist, meine Macht nahm immer zu.
\quellseite{751}

c Anlachte, freundlich, fröhlich mit ihnen war, wurden sie darum nicht kühn, mich zu verachten. Id est, familiaritas mea non peperit apud eos mei contemptum.

\section{Summarium}

Dies ist ein Klagrede, in welcher Hiob anzeigt, in was Stand und Wesen er gewesen sei, ehe er in das Unglück kommen sei. Und ist dies Kapitel ein fein Bild, weiß sich ein jeder, der da gottesfürchtig leben will, in seinem Leben befleißigen soll.

Zum andern ist auch ein Exempel, weil Hiob so fromm gewesen, und ihn dennoch Gott unter das Kreuz und Leiden geworfen hat, daß wir uns des Kreuzes auch nicht beschweren sollen, die wir doch mancherlei Ursache dazu geben. Denn, daß wir nicht allweg unser Theil nach dem, wie es äußerlich gehet, richten sollen. Sonst werden wir oft fromme Leute für böse halten und ihnen unrecht thun.

\section{Gebet}

Helfe GOtt, daß wir alle in seiner Furcht leben und wandeln, wie es eines jeden Beruf erfordert; und verleihe, daß ein jeder unter seinem Kreuz geduldig sei, und sich demüthige, bis GOtt alle aus zeitlicher Leibes- und der Seelen-Noth väterlich errettet und helfe, Amen.

\chapter[30]{30}

\section{Vorrede}

- Geliebte Christen! Es ist eine schwere Versuchung, wenn eines verneinet: Gott erhöre sein Gebet nicht. Denn er würde ihm sonst das benehmen, darum er ihn lange Zeit Tag und Nacht gebeten hat. Aber da soll man wissen, daß Gott alle, die aus wahrem Glauben zu ihm rufen, erhöre. Denn wir haben dessen vielfältige Verheißungen, und auch tröstliche Exempel. Daß er aber dergleichen thut, als ob er uns nicht erhöre, und uns gerade nicht bewilligt, daß uns an Leib und Seele beschwerlich ist; oder uns nicht gibt, daß wir von ihm begehren, das geschieht darum, daß er unsern Glauben, Hoffnung und Geduld üben will; und daß wir auf ihn warten und ihm nichts vorschreiben sollen: sondern gewiß sein, daß er zu rechter Zeit kommen und in der Weile uns nichts versäumen und verderben werde. Denn wie ein Goldschmied, der das Silber in das Feuer setzt, dabei stehet und nicht, so bald es anhebt zu schmelzen, aufhöret, am Blasebalg zu ziehen, sondern weiß, wenn es genug ist, und siehet, daß das Feuer nicht zu groß, noch zu klein sei, sondern eben recht; also ist Gott auch bei allen unsern Trübsalen, und thut allen Dingen recht zu seiner Zeit, und wenn es nur zu seinen Ehren gelanget und den Menschen nützlich ist. Dies erinnert auch das 30. Kapitel.

\section{Hiob klaget über sein Elend und Anfechtung}

a Sie die Chaldäer.

b Id est, deposuerunt, privarunt curru et aurigata, id est, dominio meo. Das heißt: abgesetzt, des Wagens und der Zügel, d. i. meiner Herrschaft beraubt.
\quellseite{752}

c) Gekleidet) das ist, mancherlei Unglück wird mir angethan gewaltiglich, daß ich mich nicht erwehren kann und gürtet mich, daß ich nicht heraus kommen kann, und muß anhaben, wie den Rock am Halse.

d) Beinhaus) das ist, im Beinhause werde ich Ruhe haben.

\section{Summarium}

Sie ist eine Klagrede, in welcher Hiob, nachdem er erzählt hat, wie wohl und glücklich es ihm vorher gegangen und wie fromm er gewesen sei, klagt, wie es jetzt alles mit einander umgekehret sei, und er solches gottfürchtigen Lebens nicht genieße. Solches, spricht Hiob, widerfähret mir darum, daß Gott nicht mehr mein Freund ist, sondern handelt gegen mir, wie ein Grausamer, wider welches Gewalt ich mich nicht kann schüßen, V. 21. Also sehen wir, wie die Vernunft pfleget zu urtheilen, wenns übel gehet, Noth und Kümmerniß für der Hand und Gottes Wort hinweg ist. Da denkt man alsobald, Gott wolle unser nicht und sei zornig über uns. So doch Gott alles leiden darum zuschicket, daß der alte Adam getödtet, der Glaube, Gehorsam und Geduld in uns probiret und gemehret, und wir also endlich selig werden.

\section{Votum}

Helfe Gott, von dem alles Vermögen kommt, Gutes zu thun, daß wir unter dem Kreuz unsere Seelen mit Geduld fassen; und verleihe, daß wir darunter allzeit als wahre Christen erfunden werden, fromm sein und auch an Leib und Seele, Amt, Nahrung, Glück und Gedeihen haben, Amen.

\chapter[31]{31}

\section{Vorrede}

Geliebte Christen! Nach dem Glauben, der die fürnehmste Tugend ist eines wahren Christen, soll ein jeder frommer Hausvater fürnehmlich sechserlei Tugenden an sich haben. Erstlich soll er sein ein keuscher Ehemann, der ihm an seinem Ehegemahl genügen läßt. Zum andern soll er sein ein Liebhaber der Wahrheit. Denn er soll nicht mit Lügen und Trügen umgehen, daß er seinen zeitlichen Nutzen schaffe, sondern er soll aufrichtig und redlich in seinem Beruf handeln. Zum dritten soll er gerecht sein gegen sein Hausgesinde, und dem erzeigen, was recht und billig ist. Zum vierten soll er gegen den Armen wohlthätig sein, und den Wittwen und Waisen mit Rath und That dienen. Zum fünften soll ein Hausvater demüthig sein in Glück und Wohlfahrt. Zum sechsten und letzten soll er gern beherbergen die Fremdlinge und zur Unschuld verjagte Leute. Das ist aller christlichen Hausväter Haustafel, darnach sie sich halten sollen. Gott, der es ihnen befohlen, zum Gehorsam: Dem heiligen Evangelio, dem sie folgen sollen; zu Ehren: Ihrem Nächsten zu Nutz und ihnen selbst zu gutem Gewissen, Nutz und Wohlfahrt, durch Gottes milden Segen. Und ein solcher Hausvater ist gewesen Hiob, von dem das 31. Kapitel.

\section{Hiobs unsträflicher Wandel}
\quellseite{753}

a Das ist, wenn mirs glücklich ginge, habe ich nicht meine Freude darinnen gehabt.

b Hand küssen: heißt sein eigen Werth preisen, welches allein Gott zugehöret.

c Das ist, mein Gesinde mußte auch nichts begehren an meine Feinde.

d Fürst) frei, unerschrocken.
\quellseite{754}

\section{Summarium}

Ist eine feine Lehre, in welcher Hiob angezeiget, daß er sein Lebenlang in Gottesfurcht und gutem Gewissen geblieben sei, B. 1. Er sei kein Hurer noch Ehebrecher gewesen; habe mit seinen Unterthanen nicht tyrannisch gefahren, B. 5. 13. Sei nicht geizig noch unbarmherzig gewesen, habe mit den Waisen nicht untreulich gehandelt, B. 16. Habe sich Geldes und Gutes nicht übernommen, B. 24. Sei keines Wanckels halben nicht hoffärtig gewesen, B. 25., habe sich ander Leute Unglück nicht gefreuet, sei gastfrei gewesen, habe niemals weder Gunst noch Ungunst angesehen in Händeln, und seinem Gesinde treulich und recht gelohnet, B. 32. 39. Solch Exempel sollen wir wohl bedenken und uns auch darnach halten. Sonderlich aber ist hie zu merken, daß er den Ehebruch ein solch Feuer heißt, welches bis ins Verderben verzehret und alles Einkommen auswurzelet, B. 12. Solches Kapitel sollten alle Hausväter in ihren Häusern haben, und täglich lesen, und ihr Leben darnach richten.

\section{Psalm}

\chapter[32]{32}

\section{Vorrede}

Bei diesem Kapitel sollen wir merken folgenden Unterricht. Die, so Gott mit Kreuz nach seinem gnädigen Willen heimsucht, soll man freilich ermahnen, warnen und strafen, wo sie in etwas, belangende Worte oder Werke, sich nicht recht erzeigen und zu viel thun; aber dabei soll man es nicht bleiben lassen, sondern man soll sie auch lehren und vermahnen, daß sie das Kreuz in Geduld tragen, weil es eine väterliche Züchtigung ist unser Sünde halben, und daß sie darunter mit gläubigem und fleißigem Gebet um Hülfe und Erlösung bei Gott, der ein Meister zu helfen ist, anhalten sollen. Denn soll man die lieben Kreuzträger trösten mit Gottes Wort, wenn sie sonderlich betrübt und kleinmüthig sind, weil das Kreuz so lange währet und ihnen oft grausame Leibes-Schmerzen anrichtet. Fürnehmlich aber soll man auch für die Kreuzträger herzlich beten, darum, daß es Gott kräftlich befohlen hat, und daß es ihnen sonderlich hoch vonnöthen ist, ihrer großen Leibes-Schmerzen und der Seele Anfechtungen halben, die oft zugleich zusammen schlagen und kommen. Derwegen soll man bitten, daß sie Gott in wahrem starken Glauben und christlicher Geduld erhalte, daß er ihnen zu seiner Zeit helfe, oder ihre Noth zu einem seligen Ende schicke. Zu solchem Gebet für alle, so unter dem Kreuz liegen, soll uns auch bewegen, daß sie ja unser Brüder oder Schwestern sein, und daß wir einen Vater haben, und ein Erbe alle zugleich im Himmel zu erwarten. Sollen uns derhalben anderer Christen Noth so fleißig, als unserer eigenen, annehmen. Neben dem soll aus dem Kapitel gemerket werden eine feine Kinderlehre, daß junge Leute mit Reden nicht behend sein, sondern den Alten weichen und zuhören sollen. Denn alte Leute in großen Sachen geübet sind, sind fürsichtig, wissen, was Nutz und Schaden bringen mag, und wegen ein Ding eher, denn daß sie es wagen. Darnach sind die Anfechtungen nicht so stark in ihnen, als in Jungen, darum sie auch bessern Rath geben können. Dieses lehret Sirach im 32. Kapitel B. 10. und sagt: Ein Jüngling mag auch wohl reden ein mal oder zwei, wenns ihm nöth ist, und wenn man ihn fragt, soll ers kurz machen, und sich halten, als der nicht viel wisse, und lieber schweigen, und soll sich nicht den Herren gleich achten, und wenn ein Alter redet, nicht drein maschen.

\section{Elihu zürnet über das Stillschweigen der Freunde Hiobs}
\quellseite{755}
* Ich ersticke sonst für großer Weisheit.

\section{Summarium}

Nachdem die drei Freunde Hiobs aufhören zu reden, darum, daß sie ihn von seiner Meinung nicht können bringen, tritt ein anderer auf, disputiret von neuem mit Hiob und strafet Hiob, daß er in seinen Leiden unrecht und wider Gott geredet habe, und läßt doch das Leben Hiobs ungetadelt, bleibet aber fürnehm(lich) darauf, Gott sei gerecht. Darum wer da achtet, es geschehe ihm unrecht, der sei bald im selbigen Unrecht, daß ers dafür hält, Gott thue ihm unrecht. Nimmt also für sich etliche Stücke, die er von Hiob gehöret und werth achtet, daß mans strafen soll.

Neben dem ist hie eine feine Kinderlehre, daß junge Leute mit Reden nicht behend sein, sondern den Alten weichen und zuhören sollen.

\section{Psalm}

\chapter[33]{33}

\section{Vorrede}

Geliebte Christen! Einen sichern ungebrochenen Menschen, der durch Kreuz und Noth nicht gedemüthiget ist, den kann der Teufel bringen und gebrauchen, wozu er will, auch zu den allerhäßlichsten und abscheulichsten Sünden und Lastern. Derwegen, daß die Menschen nicht sicher werden, und den Teufel sich verführen lassen, so legt ihnen Gott Kreuz auf, dadurch er sie in der Zucht hält und oft zu sündigen hindert, und, da er sie eine Zeitlang damit drücket und presset, so will er doch gnädig sein den Bußfertigen und endlich helfen. Derwegen soll ein jeder sein Kreuz in Geduld tragen, und soll in Gottes heiliger Gericht nicht reden, er nehme es gleich mit uns für, wie er wolle. Denn er niemand unrecht thut, als der da weiß, wer da recht oder unrecht hat. So ist er nicht allein weise, sondern die Weisheit selbst, nicht allein gerecht, sondern auch die Gerechtigkeit selbst. Dieses lehret auch das 33. Kapitel.

\section{Elihu vertheidiget Gottes Gerechtigkeit}
\quellseite{756}

a Bedenket) sicut homo post factum consuevit, paenitet, et cogitat mutare. Triumphator in Israel (inquit Samuel) non paenitet, nec mutat. Wie der Mensch nach der That überlegt, es bereut, und denkt zu verändern. Dem Held in Israel gereuet nicht (spricht Samuel), noch verändert er sich.

b Vorhaben) wie Abimelech, 1 Mos. 20, 3.

c Daß ihm's Gesicht vergehet, daß er weder siehet noch höret.

d Zwei oder drei mal) das ist, oftmals.

\section{Summarium}

Ist eine Anklage wider Hiob, daß er unrecht sei, darum, daß er über Gott geklagt habe, er thue ihm unrecht, V. 11. Denn Gott ist mehr denn ein Mensch; darum soll kein Mensch von Gott Ursache oder Rechenschaft fordern, warum er dies oder anders thut.

Darum setzet Elihu eine feine Lehre, daß Gott Kreuz und Leiden uns Menschen darum zuschickt, daß er uns vom Unglück wenden und unsre Seele vom Verderben und vor des Teufels Strick bewahren wolle, V. 16. Solches, spricht er, weiß ein Mensch von sich selbst in der Anfechtung nicht. Aber man lernets aus dem Wort Gottes, welches die Engel, das ist, die Prediger verkündigen. Dasselbige tröstet also, daß Gott wieder helfen und gnädig wolle sein, V. 23. Wer solches also weiß und für wahr hält, der fähet an zu beten und empfindet Trost und Freude, und bekennet, daß ihm Gott nicht zu viel gethan habe; sondern, daß er durch seine Sünde noch ärgers verdienet hätte.

\section{Votum}

Verleihe unser treuer lieber GOtt und Vater, daß alle, so das Kreuz tragen, Buße thun, geduldig sein und durch JEsum Christum, ihren Heiland, in wahrem Glauben der Hülfe sich trösten, unter dem Kreuz auch endlich Hülfe an Leib und Seele erlangen, Amen.

\chapter[34]{34}

\section{Vorrede}

Geliebte Christen! Die Ungeduld im Kreuz sucht das Herz oft vermessen an, und bewegt, daß es in Trübsal und Anfechtungen zu murren, zu fluchen, unleidlich zu werden und sich wider Gottes Willen aufzulehnen, und daß es darüber des Gebets und Danksagung vergisset; sonderlich wenn es siehet, wie es den Gottlosen so wohl gehet, und daß sie der betrübten und nothleidenden Gläubigen gleich spotten, als kenneten oder achteten Gott derselben nicht. Aber da ist zu wissen, daß die Ungeduld und das Murren unter dem Kreuz eine greuliche Sünde sei, und daß Gott die auch ernstlich strafe, als davon lehret die heilige Judith am 8. Kapitel B. 21. ihres Büchleins: Die andern, spricht sie, so die Trübsal nicht haben wollen annehmen, mit Gottesfurcht, sondern mit Ungeduld wider Gott gemurret und gelästert haben, sind von dem Verderber und durch die Schlangen umgebracht. Dieses wiederholet St. Paulus mit herzlicher Warnung
\quellseite{757}
und spricht 1 Cor. 10, 10.: Murret nicht, gleichwie jener etliche murreten, und wurden umgebracht durch den Verderber. Derwegen für Ungeduld unter dem Kreuz sich jedermann fleißig hüten soll. Dieses lehret auch das 34. Kapitel.

\section[Elihu schilt die Reden Hiobs, und fährt fort,]{Elihu schilt die Reden Hiobs, und fährt fort, GOttes Gerechtigkeit zu vertheidigen}

a Meinen Pfeilen) das sind Gottes Pfeile, die in mir stecken.

b Zu drängen) das ist, er lässet einen Tyrannen regieren, der das Volk mit Aufsehern und Schinden führet und quälet.
\quellseite{758}

\section{Summarium}

Ist die andere Anklage wider Hiob, daß er habe unrecht daran geredet, da er gesagt: Es gehe dem Frommen gleich wie dem Gottlosen; solches, spricht Elihu, ist eben so viel geredet, als wäre Gott auch gottlos und thäte unrecht. B. 12. Aber solches heißt Gott lästern; denn Gott ist ein Herr über alles, und könnte allwege Ursache genug wider die Menschen und alle Welt haben, daß ers in einen Haufen würfe; denn er siehet und weiß alles, warum gehet er auch also unter den Leuten um, und strafet die Gottlosen hin und wieder, wie sie es verdienen. Und du, Hiob, verdienst solche Strafe auch wohl mit deinem Lästern; und ist recht, daß Gott mit Strafe anhält, weil du von Murren und Lästern wider Gott nicht willst aufhören.

\section{Votum}

Gott, der ein Gott ist aller Geduld und Trostes, regiere alle betrübte Kreuzträger, daß sie ja unterm Kreuze nicht lästern, sondern seine heilige Gerichte lieben und loben, und in Geduld ferner Hülfe erwarten, durch Beistand und Kraft des heiligen Geistes, hochgelobet in Ewigkeit, Amen.

\chapter[35]{35}

\section{Vorrede}

Bei diesem Kapitel sollen wir merken zwei Stücke. Das erste ist, daß Gott der Gottlosen Gebet nicht erhöre, das ist, deren, die nicht ihre Sünde erkennen, lassen sie ihnen nicht leid sein, und gedenken nicht davon abzustehen, sondern vielmehr in Sünden fortzufahren. Davon sagt der Tert B. 12.: Sie werden schreien, und er wird sie nicht erhören, denn Gott wird das eitel nicht erhören. Also lehret Jesaias Kap. 59. B. 2.: Eure Untugend scheiden euch und euren Gott von einander, und eure Sünden verbergen das Angesicht von euch, daß ihr nicht gehöret werdet. Und Salomon in seinen Sprüchen im 28. Kapitel B. 9. sagt: Wer seine Ohren abwendet zu hören das Gesetz, des Gebet ist ein Greuel. Derwegen wer fruchtbarlich beten will, daß ihm und andern zu Nutz diene, der erkenne seine Sünde, und er trete ab von der Ungerechtigkeit, wie St. Paulus redet 2 Tim. 2, 19. Das ist, er bitte Gott zuvor, ehe denn er sein Gebet anfängt, um Vergebung derselben, so wird ihm widerfahren, davon Gott sagt 2 Chronika 7. Kapitel B. 14.: Wenn sie beten und mein Angesicht suchen, und sich von ihren bösen Wegen bekehren werden, so will ich vom Himmel hören. Das andere Stück ist, daß Gott endlich alles Uebel an allen Menschen strafe. Davon sagt alhie abermals Elihu B. 14. 15.: Es ist ein Gericht vor Gott, harre sein nur, ob sein Zorn so bald nicht heimsucht, und sich nicht annimmt, daß so viel Laster da sind. So bezeugets auch die Erfahrung, daß Gott den Verzug der Strafe pfleget mit der Schwerheit der Strafe einzubringen.

\section{Elihus fernere Rede von Gottes Gerechtigkeit}

a Gesänge) das ist, der Vögel Gesänge. Oder geistlich, daß man ihn lobet in Leid und Unfall, wie der 42. Psalm B. 9. auch sagt: Des Nachts singe ich ihm.
\quellseite{759}

\section{Summarium}

Dies ist die dritte Anklage wider Hiob, der Gott beschuldiget hat, daß er ihm seines gottseligen Wandels nicht genießen lasse. Solches, spricht Elihu, ist auch unrecht geredet. Ursache, wenn du gleich lange fromm bist, so hilft Gott deine Frömmigkeit nicht. Warum soll er dirs denn bezahlen. Die Leute hilft es, wenn du ihnen was Gutes thust. Wiederum, so du ihnen unrecht thust, so schadet Gott solches nicht; es schadet nur den Leuten, und Gott will solches Unrecht über sie gehen lassen und ihr Gebet nicht erhören, weil sie so eitel und gottlos sind, V. 8. Aber um solches willen soll man nicht sagen, daß Gott das Uebel nicht strafe, er hat sein Stündlein, daß er mit der Strafe kommen wird, dasselbige wird nicht außen bleiben, V. 15.

\section{Gebet}

Der allmächtige, barmherzige GOtt helfe uns, daß wir uns täglich zu ihm bekehren, und daß wir aus bußfertigen Herzen in wahrem Glauben unser Gebet thun, und lindere die wohlverdienten Strafen, um Jesu Christi, seines lieben Sohnes und unsers Herrn willen, Amen.

\chapter[36]{36}

\section{Vorrede}

Bei diesem Kapitel sollen wir lernen, daß Gott darum mit Kreuz, Leiden und Trübsal heimsucht, auf daß er dadurch den zukünftigen Sünden wehre. Denn wenn der Mensch ohne Kreuz und Widerwärtigkeit ist, und es ihm nach seinem Wunsch und Willen gehet, so wird er gemeiniglich fürwitzig, muthwillig und sicher, und gewinnet Lust und Liebe zu allerlei sündhaften Dingen. Daher kommen die Sprichwörter: Wenn dem Esel zu wohl ist, so gehet er aufs Eis tanzen. Und: Es kann ein Mensch nichts weniger erleiden, als gute Tage. Dergleichen dies: Es müssen starke Beine sein, die gute Tage ertragen können. Solches auch die Erfahrung und Exempel bezeugen. Die heilig und gottfürchtig lebte David, da ihn Saul, und darnach Absalon, sein Sohn, verfolgte. Aber da er Ruhe und Frieden hatte, fiel er und sündigte schändlich. Also that auch der König Hiskiä, da der Feind für der Stadt Jerusalem lag, war er ganz gottfürchtig und suchte bei dem Propheten Jesaia Trost und Rath. Dergleichen that er auch, da er todt krank lag; aber da sein Feind, der Sanherib, durch göttliche Macht weggetrieben und er wieder gesund worden war, geriet er in Ruhmredigkeit, welches auch seine Nachkömmlinge mußten entgelten. Daraus denn folget, daß der Mensch frömmer ist, wenn es ihm übel gehet, denn wenn es ihm wohlgehet. Derwegen so drückt uns Gott mit Kreuz, daß wir desto weniger sündigen sollen, und daß er fromme Kinder habe, die ihm dienen in Heiligkeit und Gerechtigkeit. Dies bezeugen auch St. Paulus, Röm. 6, 6. und saget: Dieser alter Mensch wird gekreuziget, auf daß der sündliche Leib aufhöre, daß wir hinfort der Sünde nicht dienen. Ja der Heyde Plinius jun. lehret in einer seiner Epistel, daß wir in Krankheiten (also unter anderm Kreuz) nicht so Lust zu sündhaften Dingen haben, als wenn wir frisch und gesund sein, und saget darauf: Wollte Gott, daß wir allezeit wären, wie wir sein, wenn wir krank liegen. Hieher gehört auch, daß St. Petrus 1. Epistel 4. Capitel V. 1. sagt: Wer am Fleisch leidet, der höret auf zu sündigen, daß er hinfort, was noch hinterstelliger Zeit im Fleisch ist, nicht der Menschen Lüsten, sondern dem Willen Gottes lebe.

\section{Ferner Beweis der Gerechtigkeit, Allmacht und Weisheit Gottes}
\quellseite{760}

a Id est, ab occidente in orientem. D. h. von West nach Ost.

\section{Summarium}

Nachdem Elihu an Hiob gestraft hat, daß er unrecht geredet und Gott gelästert habe, fähret er hie zum Beschluß eine neue Predigt an, von Gott, von seinen Werken und Gericht, B. 5., von seiner Gewalt und Allmächtigkeit, auf daß Hiob für solchem großen Gott lerne sich demüthigen, und ihn fürchten, und nicht weiter wider ihn rede.

Und ist sonderlich in diesem Kapitel zu merken, daß Elihu sagt: Gott sei gerecht, und thue niemand unrecht, B. 7. Daß er aber mancherlei Kreuz und Anfechtung den Leuten zuschicke, geschieht darum, daß sie ihre Untugend erkennen lernen, sich züchtigen lassen und von Unrecht sich bekehren. Denn wo uns Gott nicht durch mancherlei Anfechtung heimsucht, würden wir je länger, je ärger.

\section{Votum}

Wir danken GOtt herzlich, daß er durch das Kreuz unser Bestes und Nutzen suchet; der helfe allen lieben Kreuzträgern, und regiere sie, daß sie ihr Kreuz in christlicher Geduld tragen, und daß sie auch nach der Strafe frömmer werden, und beständig fromm bleiben und selig werden, Amen.

\chapter[37]{37}

\section{Vorrede}

Bei diesem Kapitel sollen wir merken folgende Lehre: Wenn der Mensch gleich kein Wort Gottes gehöret hätte, und nur das schöne herrliche Gebau des Himmels und der Erden, sammt allem, was darinnen ist und geschieht, fleißig anschauet, so kann ers nicht lassen, er muß sich darum bekümmern, wo es doch alles herkomme, es werde ja nicht ohngefähr also worden sein, es müsse dennoch irgend woher seinen Anfang haben; und schließt, daß etwas sei, nämlich Gott, der solches alles gemacht habe. Wenn der Mensch darnach weiter siehet, daß Himmel und Erden so lange Zeit, so fest und unverrücklich bestanden, Sterne, Sonne und Mond ihren Lauf so ordentlich unverhindert halten, daß Jahr seine unterschiedliche Ordnung und Abwechselung des Sommers und Winters, Herbsts und Lenzes hält, und daß eine jede Zeit ihre Eigenschaft mit sich bringet, als Regen, Schnee, Donner und Wind; und daß die Creaturen, Vögel, Fische, Thiere und Menschen, ein jedes in seinem Geschlecht und seiner Art sich vermehren und zusammen halten: Daß, zu dem auch die Gewächse, Bäume, Kräuter und andere Früchte, sich jährlich besamen; und daß eine solch schöne Ordnung allenthalben beständiglich, und solches alles dem Menschen zu gut, gehalten wird: So muß der Menschen Herz schließen, daß der, so dieses alles gemacht und erhält, und so lange Zeit so ordentlich regieret, ob er gleich unsichtbar ist, dennoch wahrhaftiger Gott und ewiger, allmächtiger, weiser, gütiger, treuer Gott sein müsse. Solche Erkenntniß Gottes findet man in dem großen Buch, welches heißt die geschaffene Welt, welches Buch so manches Blatt hat, so manche Creatur in der Welt ist, und steht an einem jeglichen Blatt geschrieben: Gott ist ewig und allmächtig, weise und gütig, gnädig und barmherzig; und sind alle Menschen darum erschaffen, und haben darum von Gott ihre Augen empfangen, daß sie in diesem schönen Buch lesen, und an solchem Werk den Schöpfer aller Creaturen erkennen, und dafür auch ihn ehren, ihm
\quellseite{761}
danken und nach seinem Willen gehorsamlich leben, und sich seiner Allmacht und Güte im Leben und Sterben freuen und trösten sollen. Derwegen, da St. Paulus zu Athen das Evangelium predigte, lehret er auch, daß man Gott in den Kreaturen suchen solle. In der Apost. 17, 27. und in der Epistel an die Römer im 1. Kapitel V. 20. schreibt er: Wie Gott diejenigen greulich gestraft, so ihn bei seinen Kreaturen und Werken nicht erkannt haben.

\section{Gottes Majestät wird aus dem Buch der Natur erkannt}

* Gold) das ist, helle Wetter, wie lauter Gold.

\section{Summarium}

Dies ist der andere Theil der Beschlußpredigt Elihu, in welcher er dem armen Hiob erzählet die herrlichen Werke Gottes, die wir alle Tage sehen, als schön Wetter, Donner, Schnee, Regen, Wind, Wolken und dergleichen, B. 2. Weil nun Hiob bekennet und sagen muß, daß solche Werke höher und größer sind, denn aller Menschen Vermögen können auswege bringen, will er ihn also vermahnet haben, daß er wider solchen hohen, großen, gewaltigen Gott nicht freventlich reden, und sein Werk, das er mit ihm fürhabe, nicht tadeln noch ausforschen soll, sondern sich für ihm demüthigen und ihm sein Werk und Willen gefallen lassen, und geduldig sein, so thut er recht.

\section{Votum}

Lob, Ehr, Preis und Dank sei GOtt für die Lehre von den Werken seiner Allmacht, Weisheit und Gütigkeit; der gebe seine Gnade, daß wir diese täglich anschauen, ihm ferner allezeit dafür danken, und dieselben wohl und christlich gebrauchen, Amen.

\chapter[38]{38}

\section{Vorrede}

Geliebte Christen! Der heilige Prophet Jeremias sagt in Klagliedern im 3. Kapitel V. 25.: Der Herr ist freundlich dem, der auf ihn harret, und der Seelen, die nach ihm fraget. Es ist ein köstlich Ding, geduldig sein, und auf die Hülfe des Herrn hoffen. Also spricht auch der Herr Christus Lucä 21, 18.: Ein Haar von eurem Haupt soll nicht umkommen. Fasset eure Seelen mit Geduld. Und St. Paulus sagt in der Epistel an die Ebräer am 10. Kapitel V. 36.: Geduld ist euch noth, auf daß ihr den Willen Gottes thut, und die Verheißung empfahet.
\quellseite{762}
Diese Sprüche lehren, daß wir unter dem Kreuz, als gottesfürchtige und gehorsame Kreuzträger, geduldig sein sollen. Damit, wenn wir auf den Zaum bringen wollten, Gott nicht verursachen, daß er uns ein schärfer Gebiß einlegen, und uns den Zaum recht machen und abrichten, oder gar verstoßen, und da wir uns nicht wollten abrichten lassen, in den Karren des höllischen Feuers anspannen müsse. Dies betrachten wir auch bei dem 38. Kapitel.

\section{Gott preiset selbst seine Allmacht und Weisheit}

a Das Siegel) das ist, ihr Stand und Wesen, daß sie gewiß sein wollen, abgebildet.

b Verborgene) das ist, ins Herz.

\section{Summarium}

Hie gehet eine andere Predigt an, welche Gott selbst dem Hiob thut, darin er alle Werke seiner göttlichen Kraft erzählet. Will dadurch den
\quellseite{763}
Hiob lernen, er habe unrecht gethan, daß er solchen großen Gott zur Rechenschaft gefordert, und wider sein Werk und Willen also gemurret habe. Uns zur Lehre, daß wir im Kreuz und Leiden geduldig sollen sein, und wider solchen Gott nicht murren, von dem wir alles haben, und ohne welchen nichts sein noch bestehen kann, beide im Himmel und auf Erden. Darum soll man solchen Gott fürchten, sein Wort und Befehl in Ehren haben und demselbigen folgen, und willig und gerne leiden, was er uns zuschickt.

\section{Votum}

Helffe GOtt, daß wir dieser nöthigen Lehre treulich folgen, und in allem Kreuz und Leiden geduldig, auch sonst zu unserer Besserung GOttes Wort gehorsam seyn, in Kraft des heiligen Geistes, um Jesu Christi, seines Sohnes und unsers Heilandes willen. Amen.

\chapter[39]{39}

\section{Vorrede}

Geliebte Christen! Im Buch der Richter im 2. Kapitel V. 1. lesen wir, daß ein Engel des HErrn, das ist, der Priester Pinehas gen Bochim kommen sei, und habe an den Leuten allda ihre Sünden durch Gottes Wort gestraft, und ihnen angekündigt, wo sie nicht in Zeiten Buße thun würden, daß sie Gott ernstlich und hart strafen wolle. Was hat aber diese Bußpredigt ausgerichtet? Davon meldet die Historia, und sagt V. 4.: Da der Engel des HErrn solche Worte geredet hatte zu allen Kindern Israel, hub das Volk seine Stimme auf und weinete, und opferten daselbst dem HErrn. Also soll man Gottes Wort hören, daß man es auch annehme, in einem feinen reinen Herzen bewahre, daraus sich bessere und dem Evangelio gehorsam sey. Dieses lehret auch David im 32. Psalm V. 8. und sagt: Ich will dich unterweisen und dir den Weg zeigen, den du wandeln sollst. Seyd nicht wie Rosse und Mauler, die nicht verständig sind. Das ist, man soll sich Gottes Wort lassen regieren, Gott fürchten und nach Gottes Wort ein gottselig Leben führen. Davon lehret auch Hiob mit seinem Exempel, in dem 39. Kapitel.

\section[An den Thieren leuchtet GOttes unerschöpfliche Weisheit hervor.]{An den Thieren leuchtet GOttes unerschöpfliche Weisheit hervor. Hiob bekennet seinen Fehler}

a) Preis) das ist, es ist nur desto trotziger und muthiger, und schnaubet, als rühmte sich, wo schreckliche Dinge, als Streit und Krieg vorhanden ist.
\quellseite{764}

b) Achtet nicht) das ist, es thut, als sei ihm nichts, das doch so schrecklich ist.

\section{Summarium}

Dies ist eine Predigt, in welcher Gott vieler Thiere gedenket, wie er sie erschaffen und einem jeden seine sonderliche Art und Eigenschaft gegeben habe, V. 1. Auf daß Hiob sich für ihm demüthige, und nicht beharre auf seinem Murren wider Gott, weil er solche Werke Gottes nicht allein nicht nachthun, sondern auch nicht verstehen kann. Solche Predigt nimmt Hiob an, bekennet seine Sünde und sagt: Er wolle es nimmer thun, V. 35.

\section{Votum}

Wir danken GOtt herzlich, daß er uns auch durch sein Wort von seinen Werken und Willen prediget; der verleihe, daß wir dies sein heiliges Wort also anhören, daß wir unsere Sünde erkennen, davon abstehen, und gegen ihm gottselig, gegen uns züchtig und gegen unsern Nächsten gerecht leben, und in solchen Tugenden selig sterben, Amen.

\chapter[40]{40}

\section{Vorrede}

Geliebte Christen! Der heilige Prophet Daniel sagt im 9. Kapitel V. 7.: Du HErr bist gerecht, wir aber müssen uns schämen. Also David im 119. Psalm V. 75. 137. spricht: HErr, du bist gerecht, und deine Gerichte sind recht. Dadurch werden wir gelehret, daß wir unter dem Kreuz gerne bekennen sollen, daß uns Gott nicht unrecht thue, sondern mit uns handele, wie wir es mit unsern Sünden wohl verdienet haben. Ferner, St. Petrus sagt 1. Epistel am 3. Kapitel V. 13. 15.: Wer ist, der euch schaden könnte, so ihr dem Guten nachkommet. Fürchtet euch für ihrem Trotzen nicht, und erschrecket nicht. Und unser HErr JEsus Christus spricht Matth. 10, 29.: Kauft man nicht zween Sperlinge um einen Pfennig? Noch fällt derselbigen keiner auf die Erde ohn euern Vater. Nun aber sind auch eure Haare auf dem Haupt alle gezählet. Darum fürchtet euch nicht; ihr seid besser, denn viel Sperlinge. Hiermit werden wir gelehret, daß uns der Teufel und die böse Welt nicht mehr schaden können, denn das ihnen Gott verhänget und zulässet. Von beiden prediget auch Gott unser HErr im 40. Kapitel Hiobs.

\section{Spiegel göttlicher Allmacht an dem Behemoth}

a. Behemoth) heißen alle großen ungeheuren Thiere. Wie Leviathan alle großen ungeheuren Fische. Aber darunter beschreibt er die Gewalt und Macht des Teufels und seines Gehülffes, des gottlosen Hausvaters in der Welt.
\quellseite{765}

\section{Summarium}

Hie ist eine harte Predigt wider Hiob. Denn hie sagt der Herr, was ihm an Hiob nicht gefallen hat. Du willst, spricht er, gerecht sein, und läßt dich dünken, ich thue dir unrecht; aber hie versuche dich, ob du die Werke kannst thun, welche ich thue, V. 2. Befindest du nun, daß du sie kannst thun, alsdann so nimm mich mit Recht für, und heiße dir Rechenschaft geben.

Daß darnach stehet von Behemoth, daß er so stark und groß sei und das ganze Erdreich ihm Speise trage, und er ganze Wasser auslaufe, V. 10.; ist zu verstehen insgemein von allen wilden, starken, ungeheuren Thieren auf Erden. Gleichwohl wird durch solchen Behemoth des Teufels und der Gottlosen Gewalt in der Welt heimlich angezeigt, daß dieselben alle ihre Gewalt nicht von ihnen selbst, sondern von Gott haben, so lange ers verhänget und duldet.

\section{Datum}

Helfe GOtt, daß wir allezeit seine Gerichte und Werke lieben und loben, derwegen auch uns unter seine gewaltige Hand demüthigen; und steure dem Teufel und der bösen Welt, daß sie uns mit ihrer List und Gewalt an unserm Leib, Seele und Leben nicht schaden, Amen.

\chapter[41]{41}

\section{Vorrede}

Geliebte Christen! Obwohl der Satan ein giftiger, listiger, mächtiger Geist ist, so kann er uns doch nichts thun ohne Gottes, des Herrn, Erlaubniß und Willen; so haben wir ja auch auf unserer Seite den großen, starken Held und Sieges-Fürsten Michael, unsern Herrn Christum, Dan. 12, 1. Der die Werke des Teufels zerstören will und kann. So bewahren uns auch die lieben Schutzengel, die heiligen und starken Engel, wenn wir in Gottesfurcht unsers ordentlichen Berufs warten, Ps. 34, 8. und 91, 11. So haben wir auch das liebe und kräftige Vater unser, und das heilige Geistes Schwert, Gottes Wort, Eph. 6, 17. Damit wir uns schützen und wehren sollen und können. Dieses sollen wir uns mit Freuden trösten wider den Teufel. Ob wir schon gegen ihm sind wie eine matte Fliege gegen einen starken Löwen, und sollen derwegen täglich in Gottesfurcht leben. Denn die Teufel fliehen für einem gottesfürchtigen Menschen, wie ein Mensch fleucht für einem grausamen Löwen, der Feuer ausspeiet, sagt der alte Irenäus. Dieß erinnert auch das 41. Kapitel.

\section{Beschreibung von des Leviathans Größe, Macht und Stärke}

a Leviathan) nennet er die großen Wasserfische im Meer, doch darunter beschreibt er der Welt Fürsten, den Teufel mit seinem Anhang.
\quellseite{766}
ᵇ Starken) das ist, die großen Fische fliehen für ihm; also auch für der Welt Gewalt fliehen die Mächtigen.

ᶜ Keine Gnade) so haben sie gesündiget, das ist, sie müssen gethan haben, und behalten, als arme Sünder.

ᵈ Weg) das ist, er schwimmet und lebet im Meer, wie er will, daß man seinen Weg von fern siehet.

\section{Summarium}

Ist der letzte Theil von der Predigt des HErrn, wider den armen Hiob, in welcher er, wie vor, dem Hiob anzeiget, er soll wider Gott nicht streiten, sondern sich leiden und geduldig sein. Den Leviathan nennt er die großen Wallfische im Meer. Aber darunter will er uns den Teufel und sein Reich fürmalen, wie gewaltig und groß es sei, daß keine menschliche Hülfe ihm kann widerstehen, aber Gott könne es und wolle es thun. Will also dem Hiob anzeigen, daß solch Leiden und Jammer der Teufel aus Verhängniß Gottes ihm habe zugerichtet. Darum soll er weder mit seiner Unschuld noch Frömmigkeit wider den Teufel fechten. Denn damit werde er nichts ausrichten. Er soll sich aber an Gott halten, und da um Hülfe und Gnade bitten, so werde der Leviathan weichen müssen.

\section{Votum}

Der HErr Zebaoth, unser HErr JEsus Christus, regiere uns, daß wir, gottesfürchtig leben; und helfe uns, daß wir mit unserm gläubigen Gebet täglich wider den Teufel streiten, wolle auch unser Gebet erhören und uns behüten für Schaden und Verderben Leibes und der Seelen, Amen.

\chapter[42]{42}

\section{Vorrede}

Geliebte Christen! Gott kann seinen Gläubigen, wenn sie gleich hinter sich kommen sind, wiederum beifür helfen, wenn sie sich nur recht nach Gottes Wort in den Handel schicken. Derwegen, wenn einer um sein Hab und Gut kömmt, durch allerlei und mancherlei Unfall, oder fällt in Verachtung oder schwere Anfechtung, so läßt es Gott nicht darum geschehen, daß er ihn gar ausschütten wolle, sondern etwa darum, daß er ihm das reichlicher wiederum gebe, daß er verloren hat. Wird aber gleich einem das Seine nicht zwiefach, so wird ihm doch oft so viel, als er vor hat, oder Gott ersetzt es in andern Wegen, daß er ihm Herz und Muth gibt, daß er ihm an wenigen genügen läßt. So übertrifft weit der geistliche Segen, den Gott den Seinen gibt, den leiblichen. So wird uns Gott auch unseres Leides erst recht ergötzen im künftigen Leben. In Summa, Gott hat niemand jemals verlassen, der ihm vertraut und sich gänzlich an ihn ergeben, und in der Noth aus wahrem Glauben zu ihm geschrieen, und ein gottesfürchtig Leben geführet hat. Davon höret auch an das 42. und letzte Capitel des Buchs Hiobs.

\section[Hiob wird auf rechte Buße von GOtt wiederum]{Hiob wird auf rechte Buße von GOtt wiederum glücklich gemacht}
\quellseite{767}

\section{Summarium}

Ist erstlich ein schön Exempel einer rechtschaffenen Buße, daß Hiob des HErrn Wort annimmt, seine Sünde erkennet und lässet ihm leid sein und bittet um Gnade, B. 2. Das aber der HErr Eliphas und die andern zwei strafet, und sagt: Hiob habe recht geredet, B. 7., gehet dahin, daß es wahr ist, Hiob habe solche Strafe nicht verdient. Denn der HErr sagt selbst: Er habe seine Hand ohne Ursache wider ihn ausgestrecket. Darum haben die Freunde Hiobs unrecht, daß sie sagen: Er müsse seiner Sünden halben leiden. Wiederum aber hat Hiob auch in dem unrecht, daß er wider Gott murret, und um solcher Strafe willen Gott für unrecht anklaget.

Zum andern ist hie ein sonderlich fein Exempel, wie Gott der Frommen Fürbitte ansehen und derselbigen andere wolle genießen lassen, B. 10.

\section{Votum}

Also hat eure Liebe die Historien vom frommen Hiob, in welcher wir fürnehmlich lernen sollen, daß GOtt auch die Frommen unter das Kreuz wirft, und der alte Adam GOttes Willen nicht kann geduldig tragen, er murret, und meinet, es geschehe ihm Unrecht. Aber tröstlich ists, daß GOtt solche Ungeduld vergeben, und aus aller Noth und Anfechtung endlich helfen, und uns alles reichlich ergänzen will. Ist derhalben eine nütze Lection für die, so in großem Unglück und Leiden stecken, daß sie hie lernen, wie sie sich für GOtt demüthigen und allein auf seine Gnade sehen, und keines Verdienstes sich rühmen sollen, so will er gnädig helfen, nicht allein hie zeitlich, sondern auch ewig. Das helfe uns unser lieber HErr Christus, Amen.
% Scan page 768: Blank family-register form excluded
% Scan page 769: No literary apparatus retained
% Scan page 770: Blank family-register form excluded
% Scan page 771: Blank family-register form excluded; Google scan watermark excluded
% Scan page 772: Blank family-register form excluded; Google scan watermark excluded
% Scan page 773: Blank family-register form excluded; Google scan watermark excluded
% Scan page 774: No literary apparatus retained
% Scan page 775: No literary apparatus retained
% Scan page 776: No literary apparatus retained
% Scan page 777: Google scan watermark excluded
% Scan page 778: Blank/non-literary scan matter excluded; Library ownership stamp matter excluded; Google scan watermark excluded
% Scan page 779: No literary apparatus retained
% Scan page 780: Library ownership stamp matter excluded; Google scan watermark excluded
% Scan page 781: Google scan watermark excluded
